\documentclass[11pt]{article}

    \usepackage[breakable]{tcolorbox}
    \usepackage{parskip} % Stop auto-indenting (to mimic markdown behaviour)
    

    % Basic figure setup, for now with no caption control since it's done
    % automatically by Pandoc (which extracts ![](path) syntax from Markdown).
    \usepackage{graphicx}
    % Keep aspect ratio if custom image width or height is specified
    \setkeys{Gin}{keepaspectratio}
    % Maintain compatibility with old templates. Remove in nbconvert 6.0
    \let\Oldincludegraphics\includegraphics
    % Ensure that by default, figures have no caption (until we provide a
    % proper Figure object with a Caption API and a way to capture that
    % in the conversion process - todo).
    \usepackage{caption}
    \DeclareCaptionFormat{nocaption}{}
    \captionsetup{format=nocaption,aboveskip=0pt,belowskip=0pt}

    \usepackage{float}
    \floatplacement{figure}{H} % forces figures to be placed at the correct location
    \usepackage{xcolor} % Allow colors to be defined
    \usepackage{enumerate} % Needed for markdown enumerations to work
    \usepackage{geometry} % Used to adjust the document margins
    \usepackage{amsmath} % Equations
    \usepackage{amssymb} % Equations
    \usepackage{textcomp} % defines textquotesingle
    % Hack from http://tex.stackexchange.com/a/47451/13684:
    \AtBeginDocument{%
        \def\PYZsq{\textquotesingle}% Upright quotes in Pygmentized code
    }
    \usepackage{upquote} % Upright quotes for verbatim code
    \usepackage{eurosym} % defines \euro

    \usepackage{iftex}
    \ifPDFTeX
        \usepackage[T1]{fontenc}
        \IfFileExists{alphabeta.sty}{
              \usepackage{alphabeta}
          }{
              \usepackage[mathletters]{ucs}
              \usepackage[utf8x]{inputenc}
          }
    \else
        \usepackage{fontspec}
        \usepackage{unicode-math}
    \fi

    \usepackage{fancyvrb} % verbatim replacement that allows latex
    \usepackage{grffile} % extends the file name processing of package graphics
                         % to support a larger range
    \makeatletter % fix for old versions of grffile with XeLaTeX
    \@ifpackagelater{grffile}{2019/11/01}
    {
      % Do nothing on new versions
    }
    {
      \def\Gread@@xetex#1{%
        \IfFileExists{"\Gin@base".bb}%
        {\Gread@eps{\Gin@base.bb}}%
        {\Gread@@xetex@aux#1}%
      }
    }
    \makeatother
    \usepackage[Export]{adjustbox} % Used to constrain images to a maximum size
    \adjustboxset{max size={0.9\linewidth}{0.9\paperheight}}

    % The hyperref package gives us a pdf with properly built
    % internal navigation ('pdf bookmarks' for the table of contents,
    % internal cross-reference links, web links for URLs, etc.)
    \usepackage{hyperref}
    % The default LaTeX title has an obnoxious amount of whitespace. By default,
    % titling removes some of it. It also provides customization options.
    \usepackage{titling}
    \usepackage{longtable} % longtable support required by pandoc >1.10
    \usepackage{booktabs}  % table support for pandoc > 1.12.2
    \usepackage{array}     % table support for pandoc >= 2.11.3
    \usepackage{calc}      % table minipage width calculation for pandoc >= 2.11.1
    \usepackage[inline]{enumitem} % IRkernel/repr support (it uses the enumerate* environment)
    \usepackage[normalem]{ulem} % ulem is needed to support strikethroughs (\sout)
                                % normalem makes italics be italics, not underlines
    \usepackage{soul}      % strikethrough (\st) support for pandoc >= 3.0.0
    \usepackage{mathrsfs}
    

    
    % Colors for the hyperref package
    \definecolor{urlcolor}{rgb}{0,.145,.698}
    \definecolor{linkcolor}{rgb}{.71,0.21,0.01}
    \definecolor{citecolor}{rgb}{.12,.54,.11}

    % ANSI colors
    \definecolor{ansi-black}{HTML}{3E424D}
    \definecolor{ansi-black-intense}{HTML}{282C36}
    \definecolor{ansi-red}{HTML}{E75C58}
    \definecolor{ansi-red-intense}{HTML}{B22B31}
    \definecolor{ansi-green}{HTML}{00A250}
    \definecolor{ansi-green-intense}{HTML}{007427}
    \definecolor{ansi-yellow}{HTML}{DDB62B}
    \definecolor{ansi-yellow-intense}{HTML}{B27D12}
    \definecolor{ansi-blue}{HTML}{208FFB}
    \definecolor{ansi-blue-intense}{HTML}{0065CA}
    \definecolor{ansi-magenta}{HTML}{D160C4}
    \definecolor{ansi-magenta-intense}{HTML}{A03196}
    \definecolor{ansi-cyan}{HTML}{60C6C8}
    \definecolor{ansi-cyan-intense}{HTML}{258F8F}
    \definecolor{ansi-white}{HTML}{C5C1B4}
    \definecolor{ansi-white-intense}{HTML}{A1A6B2}
    \definecolor{ansi-default-inverse-fg}{HTML}{FFFFFF}
    \definecolor{ansi-default-inverse-bg}{HTML}{000000}

    % common color for the border for error outputs.
    \definecolor{outerrorbackground}{HTML}{FFDFDF}

    % commands and environments needed by pandoc snippets
    % extracted from the output of `pandoc -s`
    \providecommand{\tightlist}{%
      \setlength{\itemsep}{0pt}\setlength{\parskip}{0pt}}
    \DefineVerbatimEnvironment{Highlighting}{Verbatim}{commandchars=\\\{\}}
    % Add ',fontsize=\small' for more characters per line
    \newenvironment{Shaded}{}{}
    \newcommand{\KeywordTok}[1]{\textcolor[rgb]{0.00,0.44,0.13}{\textbf{{#1}}}}
    \newcommand{\DataTypeTok}[1]{\textcolor[rgb]{0.56,0.13,0.00}{{#1}}}
    \newcommand{\DecValTok}[1]{\textcolor[rgb]{0.25,0.63,0.44}{{#1}}}
    \newcommand{\BaseNTok}[1]{\textcolor[rgb]{0.25,0.63,0.44}{{#1}}}
    \newcommand{\FloatTok}[1]{\textcolor[rgb]{0.25,0.63,0.44}{{#1}}}
    \newcommand{\CharTok}[1]{\textcolor[rgb]{0.25,0.44,0.63}{{#1}}}
    \newcommand{\StringTok}[1]{\textcolor[rgb]{0.25,0.44,0.63}{{#1}}}
    \newcommand{\CommentTok}[1]{\textcolor[rgb]{0.38,0.63,0.69}{\textit{{#1}}}}
    \newcommand{\OtherTok}[1]{\textcolor[rgb]{0.00,0.44,0.13}{{#1}}}
    \newcommand{\AlertTok}[1]{\textcolor[rgb]{1.00,0.00,0.00}{\textbf{{#1}}}}
    \newcommand{\FunctionTok}[1]{\textcolor[rgb]{0.02,0.16,0.49}{{#1}}}
    \newcommand{\RegionMarkerTok}[1]{{#1}}
    \newcommand{\ErrorTok}[1]{\textcolor[rgb]{1.00,0.00,0.00}{\textbf{{#1}}}}
    \newcommand{\NormalTok}[1]{{#1}}

    % Additional commands for more recent versions of Pandoc
    \newcommand{\ConstantTok}[1]{\textcolor[rgb]{0.53,0.00,0.00}{{#1}}}
    \newcommand{\SpecialCharTok}[1]{\textcolor[rgb]{0.25,0.44,0.63}{{#1}}}
    \newcommand{\VerbatimStringTok}[1]{\textcolor[rgb]{0.25,0.44,0.63}{{#1}}}
    \newcommand{\SpecialStringTok}[1]{\textcolor[rgb]{0.73,0.40,0.53}{{#1}}}
    \newcommand{\ImportTok}[1]{{#1}}
    \newcommand{\DocumentationTok}[1]{\textcolor[rgb]{0.73,0.13,0.13}{\textit{{#1}}}}
    \newcommand{\AnnotationTok}[1]{\textcolor[rgb]{0.38,0.63,0.69}{\textbf{\textit{{#1}}}}}
    \newcommand{\CommentVarTok}[1]{\textcolor[rgb]{0.38,0.63,0.69}{\textbf{\textit{{#1}}}}}
    \newcommand{\VariableTok}[1]{\textcolor[rgb]{0.10,0.09,0.49}{{#1}}}
    \newcommand{\ControlFlowTok}[1]{\textcolor[rgb]{0.00,0.44,0.13}{\textbf{{#1}}}}
    \newcommand{\OperatorTok}[1]{\textcolor[rgb]{0.40,0.40,0.40}{{#1}}}
    \newcommand{\BuiltInTok}[1]{{#1}}
    \newcommand{\ExtensionTok}[1]{{#1}}
    \newcommand{\PreprocessorTok}[1]{\textcolor[rgb]{0.74,0.48,0.00}{{#1}}}
    \newcommand{\AttributeTok}[1]{\textcolor[rgb]{0.49,0.56,0.16}{{#1}}}
    \newcommand{\InformationTok}[1]{\textcolor[rgb]{0.38,0.63,0.69}{\textbf{\textit{{#1}}}}}
    \newcommand{\WarningTok}[1]{\textcolor[rgb]{0.38,0.63,0.69}{\textbf{\textit{{#1}}}}}
    \makeatletter
    \newsavebox\pandoc@box
    \newcommand*\pandocbounded[1]{%
      \sbox\pandoc@box{#1}%
      % scaling factors for width and height
      \Gscale@div\@tempa\textheight{\dimexpr\ht\pandoc@box+\dp\pandoc@box\relax}%
      \Gscale@div\@tempb\linewidth{\wd\pandoc@box}%
      % select the smaller of both
      \ifdim\@tempb\p@<\@tempa\p@
        \let\@tempa\@tempb
      \fi
      % scaling accordingly (\@tempa < 1)
      \ifdim\@tempa\p@<\p@
        \scalebox{\@tempa}{\usebox\pandoc@box}%
      % scaling not needed, use as it is
      \else
        \usebox{\pandoc@box}%
      \fi
    }
    \makeatother

    % Define a nice break command that doesn't care if a line doesn't already
    % exist.
    \def\br{\hspace*{\fill} \\* }
    % Math Jax compatibility definitions
    \def\gt{>}
    \def\lt{<}
    \let\Oldtex\TeX
    \let\Oldlatex\LaTeX
    \renewcommand{\TeX}{\textrm{\Oldtex}}
    \renewcommand{\LaTeX}{\textrm{\Oldlatex}}
    % Document parameters
    % Document title
    \title{Ch13-multiple-lab}
    
    
    
    
    
    
    
% Pygments definitions
\makeatletter
\def\PY@reset{\let\PY@it=\relax \let\PY@bf=\relax%
    \let\PY@ul=\relax \let\PY@tc=\relax%
    \let\PY@bc=\relax \let\PY@ff=\relax}
\def\PY@tok#1{\csname PY@tok@#1\endcsname}
\def\PY@toks#1+{\ifx\relax#1\empty\else%
    \PY@tok{#1}\expandafter\PY@toks\fi}
\def\PY@do#1{\PY@bc{\PY@tc{\PY@ul{%
    \PY@it{\PY@bf{\PY@ff{#1}}}}}}}
\def\PY#1#2{\PY@reset\PY@toks#1+\relax+\PY@do{#2}}

\@namedef{PY@tok@w}{\def\PY@tc##1{\textcolor[rgb]{0.73,0.73,0.73}{##1}}}
\@namedef{PY@tok@c}{\let\PY@it=\textit\def\PY@tc##1{\textcolor[rgb]{0.24,0.48,0.48}{##1}}}
\@namedef{PY@tok@cp}{\def\PY@tc##1{\textcolor[rgb]{0.61,0.40,0.00}{##1}}}
\@namedef{PY@tok@k}{\let\PY@bf=\textbf\def\PY@tc##1{\textcolor[rgb]{0.00,0.50,0.00}{##1}}}
\@namedef{PY@tok@kp}{\def\PY@tc##1{\textcolor[rgb]{0.00,0.50,0.00}{##1}}}
\@namedef{PY@tok@kt}{\def\PY@tc##1{\textcolor[rgb]{0.69,0.00,0.25}{##1}}}
\@namedef{PY@tok@o}{\def\PY@tc##1{\textcolor[rgb]{0.40,0.40,0.40}{##1}}}
\@namedef{PY@tok@ow}{\let\PY@bf=\textbf\def\PY@tc##1{\textcolor[rgb]{0.67,0.13,1.00}{##1}}}
\@namedef{PY@tok@nb}{\def\PY@tc##1{\textcolor[rgb]{0.00,0.50,0.00}{##1}}}
\@namedef{PY@tok@nf}{\def\PY@tc##1{\textcolor[rgb]{0.00,0.00,1.00}{##1}}}
\@namedef{PY@tok@nc}{\let\PY@bf=\textbf\def\PY@tc##1{\textcolor[rgb]{0.00,0.00,1.00}{##1}}}
\@namedef{PY@tok@nn}{\let\PY@bf=\textbf\def\PY@tc##1{\textcolor[rgb]{0.00,0.00,1.00}{##1}}}
\@namedef{PY@tok@ne}{\let\PY@bf=\textbf\def\PY@tc##1{\textcolor[rgb]{0.80,0.25,0.22}{##1}}}
\@namedef{PY@tok@nv}{\def\PY@tc##1{\textcolor[rgb]{0.10,0.09,0.49}{##1}}}
\@namedef{PY@tok@no}{\def\PY@tc##1{\textcolor[rgb]{0.53,0.00,0.00}{##1}}}
\@namedef{PY@tok@nl}{\def\PY@tc##1{\textcolor[rgb]{0.46,0.46,0.00}{##1}}}
\@namedef{PY@tok@ni}{\let\PY@bf=\textbf\def\PY@tc##1{\textcolor[rgb]{0.44,0.44,0.44}{##1}}}
\@namedef{PY@tok@na}{\def\PY@tc##1{\textcolor[rgb]{0.41,0.47,0.13}{##1}}}
\@namedef{PY@tok@nt}{\let\PY@bf=\textbf\def\PY@tc##1{\textcolor[rgb]{0.00,0.50,0.00}{##1}}}
\@namedef{PY@tok@nd}{\def\PY@tc##1{\textcolor[rgb]{0.67,0.13,1.00}{##1}}}
\@namedef{PY@tok@s}{\def\PY@tc##1{\textcolor[rgb]{0.73,0.13,0.13}{##1}}}
\@namedef{PY@tok@sd}{\let\PY@it=\textit\def\PY@tc##1{\textcolor[rgb]{0.73,0.13,0.13}{##1}}}
\@namedef{PY@tok@si}{\let\PY@bf=\textbf\def\PY@tc##1{\textcolor[rgb]{0.64,0.35,0.47}{##1}}}
\@namedef{PY@tok@se}{\let\PY@bf=\textbf\def\PY@tc##1{\textcolor[rgb]{0.67,0.36,0.12}{##1}}}
\@namedef{PY@tok@sr}{\def\PY@tc##1{\textcolor[rgb]{0.64,0.35,0.47}{##1}}}
\@namedef{PY@tok@ss}{\def\PY@tc##1{\textcolor[rgb]{0.10,0.09,0.49}{##1}}}
\@namedef{PY@tok@sx}{\def\PY@tc##1{\textcolor[rgb]{0.00,0.50,0.00}{##1}}}
\@namedef{PY@tok@m}{\def\PY@tc##1{\textcolor[rgb]{0.40,0.40,0.40}{##1}}}
\@namedef{PY@tok@gh}{\let\PY@bf=\textbf\def\PY@tc##1{\textcolor[rgb]{0.00,0.00,0.50}{##1}}}
\@namedef{PY@tok@gu}{\let\PY@bf=\textbf\def\PY@tc##1{\textcolor[rgb]{0.50,0.00,0.50}{##1}}}
\@namedef{PY@tok@gd}{\def\PY@tc##1{\textcolor[rgb]{0.63,0.00,0.00}{##1}}}
\@namedef{PY@tok@gi}{\def\PY@tc##1{\textcolor[rgb]{0.00,0.52,0.00}{##1}}}
\@namedef{PY@tok@gr}{\def\PY@tc##1{\textcolor[rgb]{0.89,0.00,0.00}{##1}}}
\@namedef{PY@tok@ge}{\let\PY@it=\textit}
\@namedef{PY@tok@gs}{\let\PY@bf=\textbf}
\@namedef{PY@tok@ges}{\let\PY@bf=\textbf\let\PY@it=\textit}
\@namedef{PY@tok@gp}{\let\PY@bf=\textbf\def\PY@tc##1{\textcolor[rgb]{0.00,0.00,0.50}{##1}}}
\@namedef{PY@tok@go}{\def\PY@tc##1{\textcolor[rgb]{0.44,0.44,0.44}{##1}}}
\@namedef{PY@tok@gt}{\def\PY@tc##1{\textcolor[rgb]{0.00,0.27,0.87}{##1}}}
\@namedef{PY@tok@err}{\def\PY@bc##1{{\setlength{\fboxsep}{\string -\fboxrule}\fcolorbox[rgb]{1.00,0.00,0.00}{1,1,1}{\strut ##1}}}}
\@namedef{PY@tok@kc}{\let\PY@bf=\textbf\def\PY@tc##1{\textcolor[rgb]{0.00,0.50,0.00}{##1}}}
\@namedef{PY@tok@kd}{\let\PY@bf=\textbf\def\PY@tc##1{\textcolor[rgb]{0.00,0.50,0.00}{##1}}}
\@namedef{PY@tok@kn}{\let\PY@bf=\textbf\def\PY@tc##1{\textcolor[rgb]{0.00,0.50,0.00}{##1}}}
\@namedef{PY@tok@kr}{\let\PY@bf=\textbf\def\PY@tc##1{\textcolor[rgb]{0.00,0.50,0.00}{##1}}}
\@namedef{PY@tok@bp}{\def\PY@tc##1{\textcolor[rgb]{0.00,0.50,0.00}{##1}}}
\@namedef{PY@tok@fm}{\def\PY@tc##1{\textcolor[rgb]{0.00,0.00,1.00}{##1}}}
\@namedef{PY@tok@vc}{\def\PY@tc##1{\textcolor[rgb]{0.10,0.09,0.49}{##1}}}
\@namedef{PY@tok@vg}{\def\PY@tc##1{\textcolor[rgb]{0.10,0.09,0.49}{##1}}}
\@namedef{PY@tok@vi}{\def\PY@tc##1{\textcolor[rgb]{0.10,0.09,0.49}{##1}}}
\@namedef{PY@tok@vm}{\def\PY@tc##1{\textcolor[rgb]{0.10,0.09,0.49}{##1}}}
\@namedef{PY@tok@sa}{\def\PY@tc##1{\textcolor[rgb]{0.73,0.13,0.13}{##1}}}
\@namedef{PY@tok@sb}{\def\PY@tc##1{\textcolor[rgb]{0.73,0.13,0.13}{##1}}}
\@namedef{PY@tok@sc}{\def\PY@tc##1{\textcolor[rgb]{0.73,0.13,0.13}{##1}}}
\@namedef{PY@tok@dl}{\def\PY@tc##1{\textcolor[rgb]{0.73,0.13,0.13}{##1}}}
\@namedef{PY@tok@s2}{\def\PY@tc##1{\textcolor[rgb]{0.73,0.13,0.13}{##1}}}
\@namedef{PY@tok@sh}{\def\PY@tc##1{\textcolor[rgb]{0.73,0.13,0.13}{##1}}}
\@namedef{PY@tok@s1}{\def\PY@tc##1{\textcolor[rgb]{0.73,0.13,0.13}{##1}}}
\@namedef{PY@tok@mb}{\def\PY@tc##1{\textcolor[rgb]{0.40,0.40,0.40}{##1}}}
\@namedef{PY@tok@mf}{\def\PY@tc##1{\textcolor[rgb]{0.40,0.40,0.40}{##1}}}
\@namedef{PY@tok@mh}{\def\PY@tc##1{\textcolor[rgb]{0.40,0.40,0.40}{##1}}}
\@namedef{PY@tok@mi}{\def\PY@tc##1{\textcolor[rgb]{0.40,0.40,0.40}{##1}}}
\@namedef{PY@tok@il}{\def\PY@tc##1{\textcolor[rgb]{0.40,0.40,0.40}{##1}}}
\@namedef{PY@tok@mo}{\def\PY@tc##1{\textcolor[rgb]{0.40,0.40,0.40}{##1}}}
\@namedef{PY@tok@ch}{\let\PY@it=\textit\def\PY@tc##1{\textcolor[rgb]{0.24,0.48,0.48}{##1}}}
\@namedef{PY@tok@cm}{\let\PY@it=\textit\def\PY@tc##1{\textcolor[rgb]{0.24,0.48,0.48}{##1}}}
\@namedef{PY@tok@cpf}{\let\PY@it=\textit\def\PY@tc##1{\textcolor[rgb]{0.24,0.48,0.48}{##1}}}
\@namedef{PY@tok@c1}{\let\PY@it=\textit\def\PY@tc##1{\textcolor[rgb]{0.24,0.48,0.48}{##1}}}
\@namedef{PY@tok@cs}{\let\PY@it=\textit\def\PY@tc##1{\textcolor[rgb]{0.24,0.48,0.48}{##1}}}

\def\PYZbs{\char`\\}
\def\PYZus{\char`\_}
\def\PYZob{\char`\{}
\def\PYZcb{\char`\}}
\def\PYZca{\char`\^}
\def\PYZam{\char`\&}
\def\PYZlt{\char`\<}
\def\PYZgt{\char`\>}
\def\PYZsh{\char`\#}
\def\PYZpc{\char`\%}
\def\PYZdl{\char`\$}
\def\PYZhy{\char`\-}
\def\PYZsq{\char`\'}
\def\PYZdq{\char`\"}
\def\PYZti{\char`\~}
% for compatibility with earlier versions
\def\PYZat{@}
\def\PYZlb{[}
\def\PYZrb{]}
\makeatother


    % For linebreaks inside Verbatim environment from package fancyvrb.
    \makeatletter
        \newbox\Wrappedcontinuationbox
        \newbox\Wrappedvisiblespacebox
        \newcommand*\Wrappedvisiblespace {\textcolor{red}{\textvisiblespace}}
        \newcommand*\Wrappedcontinuationsymbol {\textcolor{red}{\llap{\tiny$\m@th\hookrightarrow$}}}
        \newcommand*\Wrappedcontinuationindent {3ex }
        \newcommand*\Wrappedafterbreak {\kern\Wrappedcontinuationindent\copy\Wrappedcontinuationbox}
        % Take advantage of the already applied Pygments mark-up to insert
        % potential linebreaks for TeX processing.
        %        {, <, #, %, $, ' and ": go to next line.
        %        _, }, ^, &, >, - and ~: stay at end of broken line.
        % Use of \textquotesingle for straight quote.
        \newcommand*\Wrappedbreaksatspecials {%
            \def\PYGZus{\discretionary{\char`\_}{\Wrappedafterbreak}{\char`\_}}%
            \def\PYGZob{\discretionary{}{\Wrappedafterbreak\char`\{}{\char`\{}}%
            \def\PYGZcb{\discretionary{\char`\}}{\Wrappedafterbreak}{\char`\}}}%
            \def\PYGZca{\discretionary{\char`\^}{\Wrappedafterbreak}{\char`\^}}%
            \def\PYGZam{\discretionary{\char`\&}{\Wrappedafterbreak}{\char`\&}}%
            \def\PYGZlt{\discretionary{}{\Wrappedafterbreak\char`\<}{\char`\<}}%
            \def\PYGZgt{\discretionary{\char`\>}{\Wrappedafterbreak}{\char`\>}}%
            \def\PYGZsh{\discretionary{}{\Wrappedafterbreak\char`\#}{\char`\#}}%
            \def\PYGZpc{\discretionary{}{\Wrappedafterbreak\char`\%}{\char`\%}}%
            \def\PYGZdl{\discretionary{}{\Wrappedafterbreak\char`\$}{\char`\$}}%
            \def\PYGZhy{\discretionary{\char`\-}{\Wrappedafterbreak}{\char`\-}}%
            \def\PYGZsq{\discretionary{}{\Wrappedafterbreak\textquotesingle}{\textquotesingle}}%
            \def\PYGZdq{\discretionary{}{\Wrappedafterbreak\char`\"}{\char`\"}}%
            \def\PYGZti{\discretionary{\char`\~}{\Wrappedafterbreak}{\char`\~}}%
        }
        % Some characters . , ; ? ! / are not pygmentized.
        % This macro makes them "active" and they will insert potential linebreaks
        \newcommand*\Wrappedbreaksatpunct {%
            \lccode`\~`\.\lowercase{\def~}{\discretionary{\hbox{\char`\.}}{\Wrappedafterbreak}{\hbox{\char`\.}}}%
            \lccode`\~`\,\lowercase{\def~}{\discretionary{\hbox{\char`\,}}{\Wrappedafterbreak}{\hbox{\char`\,}}}%
            \lccode`\~`\;\lowercase{\def~}{\discretionary{\hbox{\char`\;}}{\Wrappedafterbreak}{\hbox{\char`\;}}}%
            \lccode`\~`\:\lowercase{\def~}{\discretionary{\hbox{\char`\:}}{\Wrappedafterbreak}{\hbox{\char`\:}}}%
            \lccode`\~`\?\lowercase{\def~}{\discretionary{\hbox{\char`\?}}{\Wrappedafterbreak}{\hbox{\char`\?}}}%
            \lccode`\~`\!\lowercase{\def~}{\discretionary{\hbox{\char`\!}}{\Wrappedafterbreak}{\hbox{\char`\!}}}%
            \lccode`\~`\/\lowercase{\def~}{\discretionary{\hbox{\char`\/}}{\Wrappedafterbreak}{\hbox{\char`\/}}}%
            \catcode`\.\active
            \catcode`\,\active
            \catcode`\;\active
            \catcode`\:\active
            \catcode`\?\active
            \catcode`\!\active
            \catcode`\/\active
            \lccode`\~`\~
        }
    \makeatother

    \let\OriginalVerbatim=\Verbatim
    \makeatletter
    \renewcommand{\Verbatim}[1][1]{%
        %\parskip\z@skip
        \sbox\Wrappedcontinuationbox {\Wrappedcontinuationsymbol}%
        \sbox\Wrappedvisiblespacebox {\FV@SetupFont\Wrappedvisiblespace}%
        \def\FancyVerbFormatLine ##1{\hsize\linewidth
            \vtop{\raggedright\hyphenpenalty\z@\exhyphenpenalty\z@
                \doublehyphendemerits\z@\finalhyphendemerits\z@
                \strut ##1\strut}%
        }%
        % If the linebreak is at a space, the latter will be displayed as visible
        % space at end of first line, and a continuation symbol starts next line.
        % Stretch/shrink are however usually zero for typewriter font.
        \def\FV@Space {%
            \nobreak\hskip\z@ plus\fontdimen3\font minus\fontdimen4\font
            \discretionary{\copy\Wrappedvisiblespacebox}{\Wrappedafterbreak}
            {\kern\fontdimen2\font}%
        }%

        % Allow breaks at special characters using \PYG... macros.
        \Wrappedbreaksatspecials
        % Breaks at punctuation characters . , ; ? ! and / need catcode=\active
        \OriginalVerbatim[#1,codes*=\Wrappedbreaksatpunct]%
    }
    \makeatother

    % Exact colors from NB
    \definecolor{incolor}{HTML}{303F9F}
    \definecolor{outcolor}{HTML}{D84315}
    \definecolor{cellborder}{HTML}{CFCFCF}
    \definecolor{cellbackground}{HTML}{F7F7F7}

    % prompt
    \makeatletter
    \newcommand{\boxspacing}{\kern\kvtcb@left@rule\kern\kvtcb@boxsep}
    \makeatother
    \newcommand{\prompt}[4]{
        {\ttfamily\llap{{\color{#2}[#3]:\hspace{3pt}#4}}\vspace{-\baselineskip}}
    }
    

    
    % Prevent overflowing lines due to hard-to-break entities
    \sloppy
    % Setup hyperref package
    \hypersetup{
      breaklinks=true,  % so long urls are correctly broken across lines
      colorlinks=true,
      urlcolor=urlcolor,
      linkcolor=linkcolor,
      citecolor=citecolor,
      }
    % Slightly bigger margins than the latex defaults
    
    \geometry{verbose,tmargin=1in,bmargin=1in,lmargin=1in,rmargin=1in}
    
    

\begin{document}
    
    \maketitle
    
    

    
    \hypertarget{multiple-testing}{%
\section{Multiple Testing}\label{multiple-testing}}

    We include our usual imports seen in earlier labs.

    \begin{tcolorbox}[breakable, size=fbox, boxrule=1pt, pad at break*=1mm,colback=cellbackground, colframe=cellborder]
\prompt{In}{incolor}{1}{\boxspacing}
\begin{Verbatim}[commandchars=\\\{\}]
\PY{k+kn}{import}\PY{+w}{ }\PY{n+nn}{numpy}\PY{+w}{ }\PY{k}{as}\PY{+w}{ }\PY{n+nn}{np}
\PY{k+kn}{import}\PY{+w}{ }\PY{n+nn}{pandas}\PY{+w}{ }\PY{k}{as}\PY{+w}{ }\PY{n+nn}{pd}
\PY{k+kn}{import}\PY{+w}{ }\PY{n+nn}{matplotlib}\PY{n+nn}{.}\PY{n+nn}{pyplot}\PY{+w}{ }\PY{k}{as}\PY{+w}{ }\PY{n+nn}{plt}
\PY{k+kn}{import}\PY{+w}{ }\PY{n+nn}{statsmodels}\PY{n+nn}{.}\PY{n+nn}{api}\PY{+w}{ }\PY{k}{as}\PY{+w}{ }\PY{n+nn}{sm}
\PY{k+kn}{from}\PY{+w}{ }\PY{n+nn}{ISLP}\PY{+w}{ }\PY{k+kn}{import} \PY{n}{load\PYZus{}data}
\end{Verbatim}
\end{tcolorbox}

    We also collect the new imports needed for this lab.

    \begin{tcolorbox}[breakable, size=fbox, boxrule=1pt, pad at break*=1mm,colback=cellbackground, colframe=cellborder]
\prompt{In}{incolor}{2}{\boxspacing}
\begin{Verbatim}[commandchars=\\\{\}]
\PY{k+kn}{from}\PY{+w}{ }\PY{n+nn}{scipy}\PY{n+nn}{.}\PY{n+nn}{stats}\PY{+w}{ }\PY{k+kn}{import} \PYZbs{}
    \PY{p}{(}\PY{n}{ttest\PYZus{}1samp}\PY{p}{,}
     \PY{n}{ttest\PYZus{}rel}\PY{p}{,}
     \PY{n}{ttest\PYZus{}ind}\PY{p}{,}
     \PY{n}{t} \PY{k}{as} \PY{n}{t\PYZus{}dbn}\PY{p}{)}
\PY{k+kn}{from}\PY{+w}{ }\PY{n+nn}{statsmodels}\PY{n+nn}{.}\PY{n+nn}{stats}\PY{n+nn}{.}\PY{n+nn}{multicomp}\PY{+w}{ }\PY{k+kn}{import} \PYZbs{}
     \PY{n}{pairwise\PYZus{}tukeyhsd}
\PY{k+kn}{from}\PY{+w}{ }\PY{n+nn}{statsmodels}\PY{n+nn}{.}\PY{n+nn}{stats}\PY{n+nn}{.}\PY{n+nn}{multitest}\PY{+w}{ }\PY{k+kn}{import} \PYZbs{}
     \PY{n}{multipletests} \PY{k}{as} \PY{n}{mult\PYZus{}test}
\end{Verbatim}
\end{tcolorbox}

    \hypertarget{review-of-hypothesis-tests}{%
\subsection{Review of Hypothesis
Tests}\label{review-of-hypothesis-tests}}

We begin by performing some one-sample \(t\)-tests.

First we create 100 variables, each consisting of 10 observations. The
first 50 variables have mean \(0.5\) and variance \(1\), while the
others have mean \(0\) and variance \(1\).

    \begin{tcolorbox}[breakable, size=fbox, boxrule=1pt, pad at break*=1mm,colback=cellbackground, colframe=cellborder]
\prompt{In}{incolor}{3}{\boxspacing}
\begin{Verbatim}[commandchars=\\\{\}]
\PY{n}{rng} \PY{o}{=} \PY{n}{np}\PY{o}{.}\PY{n}{random}\PY{o}{.}\PY{n}{default\PYZus{}rng}\PY{p}{(}\PY{l+m+mi}{12}\PY{p}{)}
\PY{n}{X} \PY{o}{=} \PY{n}{rng}\PY{o}{.}\PY{n}{standard\PYZus{}normal}\PY{p}{(}\PY{p}{(}\PY{l+m+mi}{10}\PY{p}{,} \PY{l+m+mi}{100}\PY{p}{)}\PY{p}{)}
\PY{n}{true\PYZus{}mean} \PY{o}{=} \PY{n}{np}\PY{o}{.}\PY{n}{array}\PY{p}{(}\PY{p}{[}\PY{l+m+mf}{0.5}\PY{p}{]}\PY{o}{*}\PY{l+m+mi}{50} \PY{o}{+} \PY{p}{[}\PY{l+m+mi}{0}\PY{p}{]}\PY{o}{*}\PY{l+m+mi}{50}\PY{p}{)}
\PY{n}{X} \PY{o}{+}\PY{o}{=} \PY{n}{true\PYZus{}mean}\PY{p}{[}\PY{k+kc}{None}\PY{p}{,}\PY{p}{:}\PY{p}{]}
\end{Verbatim}
\end{tcolorbox}

    To begin, we use \texttt{ttest\_1samp()} from the \texttt{scipy.stats}
module to test \(H_{0}: \mu_1=0\), the null hypothesis that the first
variable has mean zero.

    \begin{tcolorbox}[breakable, size=fbox, boxrule=1pt, pad at break*=1mm,colback=cellbackground, colframe=cellborder]
\prompt{In}{incolor}{4}{\boxspacing}
\begin{Verbatim}[commandchars=\\\{\}]
\PY{n}{result} \PY{o}{=} \PY{n}{ttest\PYZus{}1samp}\PY{p}{(}\PY{n}{X}\PY{p}{[}\PY{p}{:}\PY{p}{,}\PY{l+m+mi}{0}\PY{p}{]}\PY{p}{,} \PY{l+m+mi}{0}\PY{p}{)}
\PY{n}{result}\PY{o}{.}\PY{n}{pvalue}
\end{Verbatim}
\end{tcolorbox}

            \begin{tcolorbox}[breakable, size=fbox, boxrule=.5pt, pad at break*=1mm, opacityfill=0]
\prompt{Out}{outcolor}{4}{\boxspacing}
\begin{Verbatim}[commandchars=\\\{\}]
0.9307442156164141
\end{Verbatim}
\end{tcolorbox}
        
    The \(p\)-value comes out to 0.931, which is not low enough to reject
the null hypothesis at level \(\alpha=0.05\). In this case,
\(\mu_1=0.5\), so the null hypothesis is false. Therefore, we have made
a Type II error by failing to reject the null hypothesis when the null
hypothesis is false.

We now test \(H_{0,j}: \mu_j=0\) for \(j=1,\ldots,100\). We compute the
100 \(p\)-values, and then construct a vector recording whether the
\(j\)th \(p\)-value is less than or equal to 0.05, in which case we
reject \(H_{0j}\), or greater than 0.05, in which case we do not reject
\(H_{0j}\), for \(j=1,\ldots,100\).

    \begin{tcolorbox}[breakable, size=fbox, boxrule=1pt, pad at break*=1mm,colback=cellbackground, colframe=cellborder]
\prompt{In}{incolor}{5}{\boxspacing}
\begin{Verbatim}[commandchars=\\\{\}]
\PY{n}{p\PYZus{}values} \PY{o}{=} \PY{n}{np}\PY{o}{.}\PY{n}{empty}\PY{p}{(}\PY{l+m+mi}{100}\PY{p}{)}
\PY{k}{for} \PY{n}{i} \PY{o+ow}{in} \PY{n+nb}{range}\PY{p}{(}\PY{l+m+mi}{100}\PY{p}{)}\PY{p}{:}
   \PY{n}{p\PYZus{}values}\PY{p}{[}\PY{n}{i}\PY{p}{]} \PY{o}{=} \PY{n}{ttest\PYZus{}1samp}\PY{p}{(}\PY{n}{X}\PY{p}{[}\PY{p}{:}\PY{p}{,}\PY{n}{i}\PY{p}{]}\PY{p}{,} \PY{l+m+mi}{0}\PY{p}{)}\PY{o}{.}\PY{n}{pvalue}
\PY{n}{decision} \PY{o}{=} \PY{n}{pd}\PY{o}{.}\PY{n}{cut}\PY{p}{(}\PY{n}{p\PYZus{}values}\PY{p}{,}
                  \PY{p}{[}\PY{l+m+mi}{0}\PY{p}{,} \PY{l+m+mf}{0.05}\PY{p}{,} \PY{l+m+mi}{1}\PY{p}{]}\PY{p}{,}
                  \PY{n}{labels}\PY{o}{=}\PY{p}{[}\PY{l+s+s1}{\PYZsq{}}\PY{l+s+s1}{Reject H0}\PY{l+s+s1}{\PYZsq{}}\PY{p}{,}
                          \PY{l+s+s1}{\PYZsq{}}\PY{l+s+s1}{Do not reject H0}\PY{l+s+s1}{\PYZsq{}}\PY{p}{]}\PY{p}{)}
\PY{n}{truth} \PY{o}{=} \PY{n}{pd}\PY{o}{.}\PY{n}{Categorical}\PY{p}{(}\PY{n}{true\PYZus{}mean} \PY{o}{==} \PY{l+m+mi}{0}\PY{p}{,}
                       \PY{n}{categories}\PY{o}{=}\PY{p}{[}\PY{k+kc}{True}\PY{p}{,} \PY{k+kc}{False}\PY{p}{]}\PY{p}{,}
                       \PY{n}{ordered}\PY{o}{=}\PY{k+kc}{True}\PY{p}{)}
\end{Verbatim}
\end{tcolorbox}

    Since this is a simulated data set, we can create a \(2 \times 2\) table
similar to Table 13.2.

    \begin{tcolorbox}[breakable, size=fbox, boxrule=1pt, pad at break*=1mm,colback=cellbackground, colframe=cellborder]
\prompt{In}{incolor}{6}{\boxspacing}
\begin{Verbatim}[commandchars=\\\{\}]
\PY{n}{pd}\PY{o}{.}\PY{n}{crosstab}\PY{p}{(}\PY{n}{decision}\PY{p}{,}
            \PY{n}{truth}\PY{p}{,}
     \PY{n}{rownames}\PY{o}{=}\PY{p}{[}\PY{l+s+s1}{\PYZsq{}}\PY{l+s+s1}{Decision}\PY{l+s+s1}{\PYZsq{}}\PY{p}{]}\PY{p}{,}
     \PY{n}{colnames}\PY{o}{=}\PY{p}{[}\PY{l+s+s1}{\PYZsq{}}\PY{l+s+s1}{H0}\PY{l+s+s1}{\PYZsq{}}\PY{p}{]}\PY{p}{)}
\end{Verbatim}
\end{tcolorbox}

            \begin{tcolorbox}[breakable, size=fbox, boxrule=.5pt, pad at break*=1mm, opacityfill=0]
\prompt{Out}{outcolor}{6}{\boxspacing}
\begin{Verbatim}[commandchars=\\\{\}]
H0                True  False
Decision
Reject H0            5     15
Do not reject H0    45     35
\end{Verbatim}
\end{tcolorbox}
        
    Therefore, at level \(\alpha=0.05\), we reject 15 of the 50 false null
hypotheses, and we incorrectly reject 5 of the true null hypotheses.
Using the notation from Section 13.3, we have \(V=5\), \(S=15\),
\(U=45\) and \(W=35\). We have set \(\alpha=0.05\), which means that we
expect to reject around 5\% of the true null hypotheses. This is in line
with the \(2 \times 2\) table above, which indicates that we rejected
\(V=5\) of the \(50\) true null hypotheses.

In the simulation above, for the false null hypotheses, the ratio of the
mean to the standard deviation was only \(0.5/1 = 0.5\). This amounts to
quite a weak signal, and it resulted in a high number of Type II errors.
Let's instead simulate data with a stronger signal, so that the ratio of
the mean to the standard deviation for the false null hypotheses equals
\(1\). We make only 10 Type II errors.

    \begin{tcolorbox}[breakable, size=fbox, boxrule=1pt, pad at break*=1mm,colback=cellbackground, colframe=cellborder]
\prompt{In}{incolor}{7}{\boxspacing}
\begin{Verbatim}[commandchars=\\\{\}]
\PY{n}{true\PYZus{}mean} \PY{o}{=} \PY{n}{np}\PY{o}{.}\PY{n}{array}\PY{p}{(}\PY{p}{[}\PY{l+m+mi}{1}\PY{p}{]}\PY{o}{*}\PY{l+m+mi}{50} \PY{o}{+} \PY{p}{[}\PY{l+m+mi}{0}\PY{p}{]}\PY{o}{*}\PY{l+m+mi}{50}\PY{p}{)}
\PY{n}{X} \PY{o}{=} \PY{n}{rng}\PY{o}{.}\PY{n}{standard\PYZus{}normal}\PY{p}{(}\PY{p}{(}\PY{l+m+mi}{10}\PY{p}{,} \PY{l+m+mi}{100}\PY{p}{)}\PY{p}{)}
\PY{n}{X} \PY{o}{+}\PY{o}{=} \PY{n}{true\PYZus{}mean}\PY{p}{[}\PY{k+kc}{None}\PY{p}{,}\PY{p}{:}\PY{p}{]}
\PY{k}{for} \PY{n}{i} \PY{o+ow}{in} \PY{n+nb}{range}\PY{p}{(}\PY{l+m+mi}{100}\PY{p}{)}\PY{p}{:}
   \PY{n}{p\PYZus{}values}\PY{p}{[}\PY{n}{i}\PY{p}{]} \PY{o}{=} \PY{n}{ttest\PYZus{}1samp}\PY{p}{(}\PY{n}{X}\PY{p}{[}\PY{p}{:}\PY{p}{,}\PY{n}{i}\PY{p}{]}\PY{p}{,} \PY{l+m+mi}{0}\PY{p}{)}\PY{o}{.}\PY{n}{pvalue}
\PY{n}{decision} \PY{o}{=} \PY{n}{pd}\PY{o}{.}\PY{n}{cut}\PY{p}{(}\PY{n}{p\PYZus{}values}\PY{p}{,}
                  \PY{p}{[}\PY{l+m+mi}{0}\PY{p}{,} \PY{l+m+mf}{0.05}\PY{p}{,} \PY{l+m+mi}{1}\PY{p}{]}\PY{p}{,}
                  \PY{n}{labels}\PY{o}{=}\PY{p}{[}\PY{l+s+s1}{\PYZsq{}}\PY{l+s+s1}{Reject H0}\PY{l+s+s1}{\PYZsq{}}\PY{p}{,}
                          \PY{l+s+s1}{\PYZsq{}}\PY{l+s+s1}{Do not reject H0}\PY{l+s+s1}{\PYZsq{}}\PY{p}{]}\PY{p}{)}
\PY{n}{truth} \PY{o}{=} \PY{n}{pd}\PY{o}{.}\PY{n}{Categorical}\PY{p}{(}\PY{n}{true\PYZus{}mean} \PY{o}{==} \PY{l+m+mi}{0}\PY{p}{,}
                       \PY{n}{categories}\PY{o}{=}\PY{p}{[}\PY{k+kc}{True}\PY{p}{,} \PY{k+kc}{False}\PY{p}{]}\PY{p}{,}
                       \PY{n}{ordered}\PY{o}{=}\PY{k+kc}{True}\PY{p}{)}
\PY{n}{pd}\PY{o}{.}\PY{n}{crosstab}\PY{p}{(}\PY{n}{decision}\PY{p}{,}
            \PY{n}{truth}\PY{p}{,}
            \PY{n}{rownames}\PY{o}{=}\PY{p}{[}\PY{l+s+s1}{\PYZsq{}}\PY{l+s+s1}{Decision}\PY{l+s+s1}{\PYZsq{}}\PY{p}{]}\PY{p}{,}
            \PY{n}{colnames}\PY{o}{=}\PY{p}{[}\PY{l+s+s1}{\PYZsq{}}\PY{l+s+s1}{H0}\PY{l+s+s1}{\PYZsq{}}\PY{p}{]}\PY{p}{)}
\end{Verbatim}
\end{tcolorbox}

            \begin{tcolorbox}[breakable, size=fbox, boxrule=.5pt, pad at break*=1mm, opacityfill=0]
\prompt{Out}{outcolor}{7}{\boxspacing}
\begin{Verbatim}[commandchars=\\\{\}]
H0                True  False
Decision
Reject H0            2     40
Do not reject H0    48     10
\end{Verbatim}
\end{tcolorbox}
        
    

    \hypertarget{family-wise-error-rate}{%
\subsection{Family-Wise Error Rate}\label{family-wise-error-rate}}

Recall from (13.5) that if the null hypothesis is true for each of \(m\)
independent hypothesis tests, then the FWER is equal to
\(1-(1-\alpha)^m\). We can use this expression to compute the FWER for
\(m=1,\ldots, 500\) and \(\alpha=0.05\), \(0.01\), and \(0.001\). We
plot the FWER for these values of \(\alpha\) in order to reproduce
Figure 13.2.

    \begin{tcolorbox}[breakable, size=fbox, boxrule=1pt, pad at break*=1mm,colback=cellbackground, colframe=cellborder]
\prompt{In}{incolor}{8}{\boxspacing}
\begin{Verbatim}[commandchars=\\\{\}]
\PY{n}{m} \PY{o}{=} \PY{n}{np}\PY{o}{.}\PY{n}{linspace}\PY{p}{(}\PY{l+m+mi}{1}\PY{p}{,} \PY{l+m+mi}{501}\PY{p}{)}
\PY{n}{fig}\PY{p}{,} \PY{n}{ax} \PY{o}{=} \PY{n}{plt}\PY{o}{.}\PY{n}{subplots}\PY{p}{(}\PY{p}{)}
\PY{p}{[}\PY{n}{ax}\PY{o}{.}\PY{n}{plot}\PY{p}{(}\PY{n}{m}\PY{p}{,}
         \PY{l+m+mi}{1} \PY{o}{\PYZhy{}} \PY{p}{(}\PY{l+m+mi}{1} \PY{o}{\PYZhy{}} \PY{n}{alpha}\PY{p}{)}\PY{o}{*}\PY{o}{*}\PY{n}{m}\PY{p}{,}
         \PY{n}{label}\PY{o}{=}\PY{l+s+sa}{r}\PY{l+s+s1}{\PYZsq{}}\PY{l+s+s1}{\PYZdl{}}\PY{l+s+s1}{\PYZbs{}}\PY{l+s+s1}{alpha=}\PY{l+s+si}{\PYZpc{}s}\PY{l+s+s1}{\PYZdl{}}\PY{l+s+s1}{\PYZsq{}} \PY{o}{\PYZpc{}} \PY{n+nb}{str}\PY{p}{(}\PY{n}{alpha}\PY{p}{)}\PY{p}{)}
         \PY{k}{for} \PY{n}{alpha} \PY{o+ow}{in} \PY{p}{[}\PY{l+m+mf}{0.05}\PY{p}{,} \PY{l+m+mf}{0.01}\PY{p}{,} \PY{l+m+mf}{0.001}\PY{p}{]}\PY{p}{]}
\PY{n}{ax}\PY{o}{.}\PY{n}{set\PYZus{}xscale}\PY{p}{(}\PY{l+s+s1}{\PYZsq{}}\PY{l+s+s1}{log}\PY{l+s+s1}{\PYZsq{}}\PY{p}{)}
\PY{n}{ax}\PY{o}{.}\PY{n}{set\PYZus{}xlabel}\PY{p}{(}\PY{l+s+s1}{\PYZsq{}}\PY{l+s+s1}{Number of Hypotheses}\PY{l+s+s1}{\PYZsq{}}\PY{p}{)}
\PY{n}{ax}\PY{o}{.}\PY{n}{set\PYZus{}ylabel}\PY{p}{(}\PY{l+s+s1}{\PYZsq{}}\PY{l+s+s1}{Family\PYZhy{}Wise Error Rate}\PY{l+s+s1}{\PYZsq{}}\PY{p}{)}
\PY{n}{ax}\PY{o}{.}\PY{n}{legend}\PY{p}{(}\PY{p}{)}
\PY{n}{ax}\PY{o}{.}\PY{n}{axhline}\PY{p}{(}\PY{l+m+mf}{0.05}\PY{p}{,} \PY{n}{c}\PY{o}{=}\PY{l+s+s1}{\PYZsq{}}\PY{l+s+s1}{k}\PY{l+s+s1}{\PYZsq{}}\PY{p}{,} \PY{n}{ls}\PY{o}{=}\PY{l+s+s1}{\PYZsq{}}\PY{l+s+s1}{\PYZhy{}\PYZhy{}}\PY{l+s+s1}{\PYZsq{}}\PY{p}{)}\PY{p}{;}
\end{Verbatim}
\end{tcolorbox}

    \begin{center}
    \adjustimage{max size={0.9\linewidth}{0.9\paperheight}}{artifacts/latex/Ch13-multiple-lab_files/artifacts/latex/Ch13-multiple-lab_17_0.png}
    \end{center}
    { \hspace*{\fill} \\}
    
    As discussed previously, even for moderate values of \(m\) such as
\(50\), the FWER exceeds \(0.05\) unless \(\alpha\) is set to a very low
value, such as \(0.001\). Of course, the problem with setting \(\alpha\)
to such a low value is that we are likely to make a number of Type II
errors: in other words, our power is very low.

We now conduct a one-sample \(t\)-test for each of the first five
managers in the\\
\texttt{Fund} dataset, in order to test the null hypothesis that the
\(j\)th fund manager's mean return equals zero, \(H_{0,j}: \mu_j=0\).

    \begin{tcolorbox}[breakable, size=fbox, boxrule=1pt, pad at break*=1mm,colback=cellbackground, colframe=cellborder]
\prompt{In}{incolor}{9}{\boxspacing}
\begin{Verbatim}[commandchars=\\\{\}]
\PY{n}{Fund} \PY{o}{=} \PY{n}{load\PYZus{}data}\PY{p}{(}\PY{l+s+s1}{\PYZsq{}}\PY{l+s+s1}{Fund}\PY{l+s+s1}{\PYZsq{}}\PY{p}{)}
\PY{n}{fund\PYZus{}mini} \PY{o}{=} \PY{n}{Fund}\PY{o}{.}\PY{n}{iloc}\PY{p}{[}\PY{p}{:}\PY{p}{,}\PY{p}{:}\PY{l+m+mi}{5}\PY{p}{]}
\PY{n}{fund\PYZus{}mini\PYZus{}pvals} \PY{o}{=} \PY{n}{np}\PY{o}{.}\PY{n}{empty}\PY{p}{(}\PY{l+m+mi}{5}\PY{p}{)}
\PY{k}{for} \PY{n}{i} \PY{o+ow}{in} \PY{n+nb}{range}\PY{p}{(}\PY{l+m+mi}{5}\PY{p}{)}\PY{p}{:}
    \PY{n}{fund\PYZus{}mini\PYZus{}pvals}\PY{p}{[}\PY{n}{i}\PY{p}{]} \PY{o}{=} \PY{n}{ttest\PYZus{}1samp}\PY{p}{(}\PY{n}{fund\PYZus{}mini}\PY{o}{.}\PY{n}{iloc}\PY{p}{[}\PY{p}{:}\PY{p}{,}\PY{n}{i}\PY{p}{]}\PY{p}{,} \PY{l+m+mi}{0}\PY{p}{)}\PY{o}{.}\PY{n}{pvalue}
\PY{n}{fund\PYZus{}mini\PYZus{}pvals}
\end{Verbatim}
\end{tcolorbox}

            \begin{tcolorbox}[breakable, size=fbox, boxrule=.5pt, pad at break*=1mm, opacityfill=0]
\prompt{Out}{outcolor}{9}{\boxspacing}
\begin{Verbatim}[commandchars=\\\{\}]
array([0.00620236, 0.91827115, 0.01160098, 0.6005396 , 0.75578151])
\end{Verbatim}
\end{tcolorbox}
        
    The \(p\)-values are low for Managers One and Three, and high for the
other three managers. However, we cannot simply reject \(H_{0,1}\) and
\(H_{0,3}\), since this would fail to account for the multiple testing
that we have performed. Instead, we will conduct Bonferroni's method and
Holm's method to control the FWER.

To do this, we use the \texttt{multipletests()} function from the
\texttt{statsmodels} module (abbreviated to \texttt{mult\_test()}).
Given the \(p\)-values, for methods like Holm and Bonferroni the
function outputs adjusted \(p\)-values, which can be thought of as a new
set of \(p\)-values that have been corrected for multiple testing. If
the adjusted \(p\)-value for a given hypothesis is less than or equal to
\(\alpha\), then that hypothesis can be rejected while maintaining a
FWER of no more than \(\alpha\). In other words, for such methods, the
adjusted \(p\)-values resulting from the \texttt{multipletests()}
function can simply be compared to the desired FWER in order to
determine whether or not to reject each hypothesis. We will later see
that we can use the same function to control FDR as well.

    The \texttt{mult\_test()} function takes \(p\)-values and a
\texttt{method} argument, as well as an optional \texttt{alpha}
argument. It returns the decisions (\texttt{reject} below) as well as
the adjusted \(p\)-values (\texttt{bonf}).

    \begin{tcolorbox}[breakable, size=fbox, boxrule=1pt, pad at break*=1mm,colback=cellbackground, colframe=cellborder]
\prompt{In}{incolor}{10}{\boxspacing}
\begin{Verbatim}[commandchars=\\\{\}]
\PY{n}{reject}\PY{p}{,} \PY{n}{bonf} \PY{o}{=} \PY{n}{mult\PYZus{}test}\PY{p}{(}\PY{n}{fund\PYZus{}mini\PYZus{}pvals}\PY{p}{,} \PY{n}{method} \PY{o}{=} \PY{l+s+s2}{\PYZdq{}}\PY{l+s+s2}{bonferroni}\PY{l+s+s2}{\PYZdq{}}\PY{p}{)}\PY{p}{[}\PY{p}{:}\PY{l+m+mi}{2}\PY{p}{]}
\PY{n}{reject}
\end{Verbatim}
\end{tcolorbox}

            \begin{tcolorbox}[breakable, size=fbox, boxrule=.5pt, pad at break*=1mm, opacityfill=0]
\prompt{Out}{outcolor}{10}{\boxspacing}
\begin{Verbatim}[commandchars=\\\{\}]
array([ True, False, False, False, False])
\end{Verbatim}
\end{tcolorbox}
        
    The \(p\)-values \texttt{bonf} are simply the
\texttt{fund\_mini\_pvalues} multiplied by 5 and truncated to be less
than or equal to 1.

    \begin{tcolorbox}[breakable, size=fbox, boxrule=1pt, pad at break*=1mm,colback=cellbackground, colframe=cellborder]
\prompt{In}{incolor}{11}{\boxspacing}
\begin{Verbatim}[commandchars=\\\{\}]
\PY{n}{bonf}\PY{p}{,} \PY{n}{np}\PY{o}{.}\PY{n}{minimum}\PY{p}{(}\PY{n}{fund\PYZus{}mini\PYZus{}pvals} \PY{o}{*} \PY{l+m+mi}{5}\PY{p}{,} \PY{l+m+mi}{1}\PY{p}{)}
\end{Verbatim}
\end{tcolorbox}

            \begin{tcolorbox}[breakable, size=fbox, boxrule=.5pt, pad at break*=1mm, opacityfill=0]
\prompt{Out}{outcolor}{11}{\boxspacing}
\begin{Verbatim}[commandchars=\\\{\}]
(array([0.03101178, 1.        , 0.05800491, 1.        , 1.        ]),
 array([0.03101178, 1.        , 0.05800491, 1.        , 1.        ]))
\end{Verbatim}
\end{tcolorbox}
        
    Therefore, using Bonferroni's method, we are able to reject the null
hypothesis only for Manager One while controlling FWER at \(0.05\).

By contrast, using Holm's method, the adjusted \(p\)-values indicate
that we can reject the null hypotheses for Managers One and Three at a
FWER of \(0.05\).

    \begin{tcolorbox}[breakable, size=fbox, boxrule=1pt, pad at break*=1mm,colback=cellbackground, colframe=cellborder]
\prompt{In}{incolor}{12}{\boxspacing}
\begin{Verbatim}[commandchars=\\\{\}]
\PY{n}{mult\PYZus{}test}\PY{p}{(}\PY{n}{fund\PYZus{}mini\PYZus{}pvals}\PY{p}{,} \PY{n}{method} \PY{o}{=} \PY{l+s+s2}{\PYZdq{}}\PY{l+s+s2}{holm}\PY{l+s+s2}{\PYZdq{}}\PY{p}{,} \PY{n}{alpha}\PY{o}{=}\PY{l+m+mf}{0.05}\PY{p}{)}\PY{p}{[}\PY{p}{:}\PY{l+m+mi}{2}\PY{p}{]}
\end{Verbatim}
\end{tcolorbox}

            \begin{tcolorbox}[breakable, size=fbox, boxrule=.5pt, pad at break*=1mm, opacityfill=0]
\prompt{Out}{outcolor}{12}{\boxspacing}
\begin{Verbatim}[commandchars=\\\{\}]
(array([ True, False,  True, False, False]),
 array([0.03101178, 1.        , 0.04640393, 1.        , 1.        ]))
\end{Verbatim}
\end{tcolorbox}
        
    As discussed previously, Manager One seems to perform particularly well,
whereas Manager Two has poor performance.

    \begin{tcolorbox}[breakable, size=fbox, boxrule=1pt, pad at break*=1mm,colback=cellbackground, colframe=cellborder]
\prompt{In}{incolor}{13}{\boxspacing}
\begin{Verbatim}[commandchars=\\\{\}]
\PY{n}{fund\PYZus{}mini}\PY{o}{.}\PY{n}{mean}\PY{p}{(}\PY{p}{)}
\end{Verbatim}
\end{tcolorbox}

            \begin{tcolorbox}[breakable, size=fbox, boxrule=.5pt, pad at break*=1mm, opacityfill=0]
\prompt{Out}{outcolor}{13}{\boxspacing}
\begin{Verbatim}[commandchars=\\\{\}]
Manager1    3.0
Manager2   -0.1
Manager3    2.8
Manager4    0.5
Manager5    0.3
dtype: float64
\end{Verbatim}
\end{tcolorbox}
        
    Is there evidence of a meaningful difference in performance between
these two managers? We can check this by performing a paired \(t\)-test
using the \texttt{ttest\_rel()} function from \texttt{scipy.stats}:

    \begin{tcolorbox}[breakable, size=fbox, boxrule=1pt, pad at break*=1mm,colback=cellbackground, colframe=cellborder]
\prompt{In}{incolor}{14}{\boxspacing}
\begin{Verbatim}[commandchars=\\\{\}]
\PY{n}{ttest\PYZus{}rel}\PY{p}{(}\PY{n}{fund\PYZus{}mini}\PY{p}{[}\PY{l+s+s1}{\PYZsq{}}\PY{l+s+s1}{Manager1}\PY{l+s+s1}{\PYZsq{}}\PY{p}{]}\PY{p}{,}
          \PY{n}{fund\PYZus{}mini}\PY{p}{[}\PY{l+s+s1}{\PYZsq{}}\PY{l+s+s1}{Manager2}\PY{l+s+s1}{\PYZsq{}}\PY{p}{]}\PY{p}{)}\PY{o}{.}\PY{n}{pvalue}
\end{Verbatim}
\end{tcolorbox}

            \begin{tcolorbox}[breakable, size=fbox, boxrule=.5pt, pad at break*=1mm, opacityfill=0]
\prompt{Out}{outcolor}{14}{\boxspacing}
\begin{Verbatim}[commandchars=\\\{\}]
0.038391072368079586
\end{Verbatim}
\end{tcolorbox}
        
    The test results in a \(p\)-value of 0.038, suggesting a statistically
significant difference.

However, we decided to perform this test only after examining the data
and noting that Managers One and Two had the highest and lowest mean
performances. In a sense, this means that we have implicitly performed
\({5 \choose 2} = 5(5-1)/2=10\) hypothesis tests, rather than just one,
as discussed in Section 13.3.2. Hence, we use the
\texttt{pairwise\_tukeyhsd()} function from
\texttt{statsmodels.stats.multicomp} to apply Tukey's method in order to
adjust for multiple testing. This function takes as input a fitted
\emph{ANOVA} regression model, which is essentially just a linear
regression in which all of the predictors are qualitative. In this case,
the response consists of the monthly excess returns achieved by each
manager, and the predictor indicates the manager to which each return
corresponds.

    \begin{tcolorbox}[breakable, size=fbox, boxrule=1pt, pad at break*=1mm,colback=cellbackground, colframe=cellborder]
\prompt{In}{incolor}{15}{\boxspacing}
\begin{Verbatim}[commandchars=\\\{\}]
\PY{n}{returns} \PY{o}{=} \PY{n}{np}\PY{o}{.}\PY{n}{hstack}\PY{p}{(}\PY{p}{[}\PY{n}{fund\PYZus{}mini}\PY{o}{.}\PY{n}{iloc}\PY{p}{[}\PY{p}{:}\PY{p}{,}\PY{n}{i}\PY{p}{]} \PY{k}{for} \PY{n}{i} \PY{o+ow}{in} \PY{n+nb}{range}\PY{p}{(}\PY{l+m+mi}{5}\PY{p}{)}\PY{p}{]}\PY{p}{)}
\PY{n}{managers} \PY{o}{=} \PY{n}{np}\PY{o}{.}\PY{n}{hstack}\PY{p}{(}\PY{p}{[}\PY{p}{[}\PY{n}{i}\PY{o}{+}\PY{l+m+mi}{1}\PY{p}{]}\PY{o}{*}\PY{l+m+mi}{50} \PY{k}{for} \PY{n}{i} \PY{o+ow}{in} \PY{n+nb}{range}\PY{p}{(}\PY{l+m+mi}{5}\PY{p}{)}\PY{p}{]}\PY{p}{)}
\PY{n}{tukey} \PY{o}{=} \PY{n}{pairwise\PYZus{}tukeyhsd}\PY{p}{(}\PY{n}{returns}\PY{p}{,} \PY{n}{managers}\PY{p}{)}
\PY{n+nb}{print}\PY{p}{(}\PY{n}{tukey}\PY{o}{.}\PY{n}{summary}\PY{p}{(}\PY{p}{)}\PY{p}{)}
\end{Verbatim}
\end{tcolorbox}

    \begin{Verbatim}[commandchars=\\\{\}]
Multiple Comparison of Means - Tukey HSD, FWER=0.05
===================================================
group1 group2 meandiff p-adj   lower  upper  reject
---------------------------------------------------
     1      2     -3.1 0.1862 -6.9865 0.7865  False
     1      3     -0.2 0.9999 -4.0865 3.6865  False
     1      4     -2.5 0.3948 -6.3865 1.3865  False
     1      5     -2.7 0.3152 -6.5865 1.1865  False
     2      3      2.9 0.2453 -0.9865 6.7865  False
     2      4      0.6 0.9932 -3.2865 4.4865  False
     2      5      0.4 0.9986 -3.4865 4.2865  False
     3      4     -2.3  0.482 -6.1865 1.5865  False
     3      5     -2.5 0.3948 -6.3865 1.3865  False
     4      5     -0.2 0.9999 -4.0865 3.6865  False
---------------------------------------------------
    \end{Verbatim}

    The \texttt{pairwise\_tukeyhsd()} function provides confidence intervals
for the difference between each pair of managers (\texttt{lower} and
\texttt{upper}), as well as a \(p\)-value. All of these quantities have
been adjusted for multiple testing. Notice that the \(p\)-value for the
difference between Managers One and Two has increased from \(0.038\) to
\(0.186\), so there is no longer clear evidence of a difference between
the managers' performances. We can plot the confidence intervals for the
pairwise comparisons using the \texttt{plot\_simultaneous()} method of
\texttt{tukey}. Any pair of intervals that don't overlap indicates a
significant difference at the nominal level of 0.05. In this case, no
differences are considered significant as reported in the table above.

    \begin{tcolorbox}[breakable, size=fbox, boxrule=1pt, pad at break*=1mm,colback=cellbackground, colframe=cellborder]
\prompt{In}{incolor}{16}{\boxspacing}
\begin{Verbatim}[commandchars=\\\{\}]
\PY{n}{fig}\PY{p}{,} \PY{n}{ax} \PY{o}{=} \PY{n}{plt}\PY{o}{.}\PY{n}{subplots}\PY{p}{(}\PY{n}{figsize}\PY{o}{=}\PY{p}{(}\PY{l+m+mi}{8}\PY{p}{,}\PY{l+m+mi}{8}\PY{p}{)}\PY{p}{)}
\PY{n}{tukey}\PY{o}{.}\PY{n}{plot\PYZus{}simultaneous}\PY{p}{(}\PY{n}{ax}\PY{o}{=}\PY{n}{ax}\PY{p}{)}\PY{p}{;}
\end{Verbatim}
\end{tcolorbox}

    \begin{center}
    \adjustimage{max size={0.9\linewidth}{0.9\paperheight}}{artifacts/latex/Ch13-multiple-lab_files/artifacts/latex/Ch13-multiple-lab_34_0.png}
    \end{center}
    { \hspace*{\fill} \\}
    
    \hypertarget{false-discovery-rate}{%
\subsection{False Discovery Rate}\label{false-discovery-rate}}

Now we perform hypothesis tests for all 2,000 fund managers in the
\texttt{Fund} dataset. We perform a one-sample \(t\)-test of
\(H_{0,j}: \mu_j=0\), which states that the \(j\)th fund manager's mean
return is zero.

    \begin{tcolorbox}[breakable, size=fbox, boxrule=1pt, pad at break*=1mm,colback=cellbackground, colframe=cellborder]
\prompt{In}{incolor}{17}{\boxspacing}
\begin{Verbatim}[commandchars=\\\{\}]
\PY{n}{fund\PYZus{}pvalues} \PY{o}{=} \PY{n}{np}\PY{o}{.}\PY{n}{empty}\PY{p}{(}\PY{l+m+mi}{2000}\PY{p}{)}
\PY{k}{for} \PY{n}{i}\PY{p}{,} \PY{n}{manager} \PY{o+ow}{in} \PY{n+nb}{enumerate}\PY{p}{(}\PY{n}{Fund}\PY{o}{.}\PY{n}{columns}\PY{p}{)}\PY{p}{:}
    \PY{n}{fund\PYZus{}pvalues}\PY{p}{[}\PY{n}{i}\PY{p}{]} \PY{o}{=} \PY{n}{ttest\PYZus{}1samp}\PY{p}{(}\PY{n}{Fund}\PY{p}{[}\PY{n}{manager}\PY{p}{]}\PY{p}{,} \PY{l+m+mi}{0}\PY{p}{)}\PY{o}{.}\PY{n}{pvalue}
\end{Verbatim}
\end{tcolorbox}

    There are far too many managers to consider trying to control the FWER.
Instead, we focus on controlling the FDR: that is, the expected fraction
of rejected null hypotheses that are actually false positives. The
\texttt{multipletests()} function (abbreviated \texttt{mult\_test()})
can be used to carry out the Benjamini--Hochberg procedure.

    \begin{tcolorbox}[breakable, size=fbox, boxrule=1pt, pad at break*=1mm,colback=cellbackground, colframe=cellborder]
\prompt{In}{incolor}{18}{\boxspacing}
\begin{Verbatim}[commandchars=\\\{\}]
\PY{n}{fund\PYZus{}qvalues} \PY{o}{=} \PY{n}{mult\PYZus{}test}\PY{p}{(}\PY{n}{fund\PYZus{}pvalues}\PY{p}{,} \PY{n}{method} \PY{o}{=} \PY{l+s+s2}{\PYZdq{}}\PY{l+s+s2}{fdr\PYZus{}bh}\PY{l+s+s2}{\PYZdq{}}\PY{p}{)}\PY{p}{[}\PY{l+m+mi}{1}\PY{p}{]}
\PY{n}{fund\PYZus{}qvalues}\PY{p}{[}\PY{p}{:}\PY{l+m+mi}{10}\PY{p}{]}
\end{Verbatim}
\end{tcolorbox}

            \begin{tcolorbox}[breakable, size=fbox, boxrule=.5pt, pad at break*=1mm, opacityfill=0]
\prompt{Out}{outcolor}{18}{\boxspacing}
\begin{Verbatim}[commandchars=\\\{\}]
array([0.08988921, 0.991491  , 0.12211561, 0.92342997, 0.95603587,
       0.07513802, 0.0767015 , 0.07513802, 0.07513802, 0.07513802])
\end{Verbatim}
\end{tcolorbox}
        
    The \emph{q-values} output by the Benjamini--Hochberg procedure can be
interpreted as the smallest FDR threshold at which we would reject a
particular null hypothesis. For instance, a \(q\)-value of \(0.1\)
indicates that we can reject the corresponding null hypothesis at an FDR
of 10\% or greater, but that we cannot reject the null hypothesis at an
FDR below 10\%.

If we control the FDR at 10\%, then for how many of the fund managers
can we reject \(H_{0,j}: \mu_j=0\)?

    \begin{tcolorbox}[breakable, size=fbox, boxrule=1pt, pad at break*=1mm,colback=cellbackground, colframe=cellborder]
\prompt{In}{incolor}{19}{\boxspacing}
\begin{Verbatim}[commandchars=\\\{\}]
\PY{p}{(}\PY{n}{fund\PYZus{}qvalues} \PY{o}{\PYZlt{}}\PY{o}{=} \PY{l+m+mf}{0.1}\PY{p}{)}\PY{o}{.}\PY{n}{sum}\PY{p}{(}\PY{p}{)}
\end{Verbatim}
\end{tcolorbox}

            \begin{tcolorbox}[breakable, size=fbox, boxrule=.5pt, pad at break*=1mm, opacityfill=0]
\prompt{Out}{outcolor}{19}{\boxspacing}
\begin{Verbatim}[commandchars=\\\{\}]
146
\end{Verbatim}
\end{tcolorbox}
        
    We find that 146 of the 2,000 fund managers have a \(q\)-value below
0.1; therefore, we are able to conclude that 146 of the fund managers
beat the market at an FDR of 10\%. Only about 15 (10\% of 146) of these
fund managers are likely to be false discoveries.

By contrast, if we had instead used Bonferroni's method to control the
FWER at level \(\alpha=0.1\), then we would have failed to reject any
null hypotheses!

    \begin{tcolorbox}[breakable, size=fbox, boxrule=1pt, pad at break*=1mm,colback=cellbackground, colframe=cellborder]
\prompt{In}{incolor}{20}{\boxspacing}
\begin{Verbatim}[commandchars=\\\{\}]
\PY{p}{(}\PY{n}{fund\PYZus{}pvalues} \PY{o}{\PYZlt{}}\PY{o}{=} \PY{l+m+mf}{0.1} \PY{o}{/} \PY{l+m+mi}{2000}\PY{p}{)}\PY{o}{.}\PY{n}{sum}\PY{p}{(}\PY{p}{)}
\end{Verbatim}
\end{tcolorbox}

            \begin{tcolorbox}[breakable, size=fbox, boxrule=.5pt, pad at break*=1mm, opacityfill=0]
\prompt{Out}{outcolor}{20}{\boxspacing}
\begin{Verbatim}[commandchars=\\\{\}]
0
\end{Verbatim}
\end{tcolorbox}
        
    Figure 13.6 displays the ordered \(p\)-values,
\(p_{(1)} \leq p_{(2)} \leq \cdots \leq p_{(2000)}\), for the
\texttt{Fund} dataset, as well as the threshold for rejection by the
Benjamini--Hochberg procedure. Recall that the Benjamini--Hochberg
procedure identifies the largest \(p\)-value such that \(p_{(j)}<qj/m\),
and rejects all hypotheses for which the \(p\)-value is less than or
equal to \(p_{(j)}\). In the code below, we implement the
Benjamini--Hochberg procedure ourselves, in order to illustrate how it
works. We first order the \(p\)-values. We then identify all
\(p\)-values that satisfy \(p_{(j)}<qj/m\) (\texttt{sorted\_set\_}).
Finally, \texttt{selected\_} is a boolean array indicating which
\(p\)-values are less than or equal to the largest \(p\)-value in
\texttt{sorted\_{[}sorted\_set\_{]}}. Therefore, \texttt{selected\_}
indexes the \(p\)-values rejected by the Benjamini--Hochberg procedure.

    \begin{tcolorbox}[breakable, size=fbox, boxrule=1pt, pad at break*=1mm,colback=cellbackground, colframe=cellborder]
\prompt{In}{incolor}{21}{\boxspacing}
\begin{Verbatim}[commandchars=\\\{\}]
\PY{n}{sorted\PYZus{}} \PY{o}{=} \PY{n}{np}\PY{o}{.}\PY{n}{sort}\PY{p}{(}\PY{n}{fund\PYZus{}pvalues}\PY{p}{)}
\PY{n}{m} \PY{o}{=} \PY{n}{fund\PYZus{}pvalues}\PY{o}{.}\PY{n}{shape}\PY{p}{[}\PY{l+m+mi}{0}\PY{p}{]}
\PY{n}{q} \PY{o}{=} \PY{l+m+mf}{0.1}
\PY{n}{sorted\PYZus{}set\PYZus{}} \PY{o}{=} \PY{n}{np}\PY{o}{.}\PY{n}{where}\PY{p}{(}\PY{n}{sorted\PYZus{}} \PY{o}{\PYZlt{}} \PY{n}{q} \PY{o}{*} \PY{n}{np}\PY{o}{.}\PY{n}{linspace}\PY{p}{(}\PY{l+m+mi}{1}\PY{p}{,} \PY{n}{m}\PY{p}{,} \PY{n}{m}\PY{p}{)} \PY{o}{/} \PY{n}{m}\PY{p}{)}\PY{p}{[}\PY{l+m+mi}{0}\PY{p}{]}
\PY{k}{if} \PY{n}{sorted\PYZus{}set\PYZus{}}\PY{o}{.}\PY{n}{shape}\PY{p}{[}\PY{l+m+mi}{0}\PY{p}{]} \PY{o}{\PYZgt{}} \PY{l+m+mi}{0}\PY{p}{:}
    \PY{n}{selected\PYZus{}} \PY{o}{=} \PY{n}{fund\PYZus{}pvalues} \PY{o}{\PYZlt{}} \PY{n}{sorted\PYZus{}}\PY{p}{[}\PY{n}{sorted\PYZus{}set\PYZus{}}\PY{p}{]}\PY{o}{.}\PY{n}{max}\PY{p}{(}\PY{p}{)}
    \PY{n}{sorted\PYZus{}set\PYZus{}} \PY{o}{=} \PY{n}{np}\PY{o}{.}\PY{n}{arange}\PY{p}{(}\PY{n}{sorted\PYZus{}set\PYZus{}}\PY{o}{.}\PY{n}{max}\PY{p}{(}\PY{p}{)}\PY{p}{)}
\PY{k}{else}\PY{p}{:}
    \PY{n}{selected\PYZus{}} \PY{o}{=} \PY{p}{[}\PY{p}{]}
    \PY{n}{sorted\PYZus{}set\PYZus{}} \PY{o}{=} \PY{p}{[}\PY{p}{]}
\end{Verbatim}
\end{tcolorbox}

    We now reproduce the middle panel of Figure 13.6.

    \begin{tcolorbox}[breakable, size=fbox, boxrule=1pt, pad at break*=1mm,colback=cellbackground, colframe=cellborder]
\prompt{In}{incolor}{22}{\boxspacing}
\begin{Verbatim}[commandchars=\\\{\}]
\PY{n}{fig}\PY{p}{,} \PY{n}{ax} \PY{o}{=} \PY{n}{plt}\PY{o}{.}\PY{n}{subplots}\PY{p}{(}\PY{p}{)}
\PY{n}{ax}\PY{o}{.}\PY{n}{scatter}\PY{p}{(}\PY{n}{np}\PY{o}{.}\PY{n}{arange}\PY{p}{(}\PY{l+m+mi}{0}\PY{p}{,} \PY{n}{sorted\PYZus{}}\PY{o}{.}\PY{n}{shape}\PY{p}{[}\PY{l+m+mi}{0}\PY{p}{]}\PY{p}{)} \PY{o}{+} \PY{l+m+mi}{1}\PY{p}{,}
           \PY{n}{sorted\PYZus{}}\PY{p}{,} \PY{n}{s}\PY{o}{=}\PY{l+m+mi}{10}\PY{p}{)}
\PY{n}{ax}\PY{o}{.}\PY{n}{set\PYZus{}yscale}\PY{p}{(}\PY{l+s+s1}{\PYZsq{}}\PY{l+s+s1}{log}\PY{l+s+s1}{\PYZsq{}}\PY{p}{)}
\PY{n}{ax}\PY{o}{.}\PY{n}{set\PYZus{}xscale}\PY{p}{(}\PY{l+s+s1}{\PYZsq{}}\PY{l+s+s1}{log}\PY{l+s+s1}{\PYZsq{}}\PY{p}{)}
\PY{n}{ax}\PY{o}{.}\PY{n}{set\PYZus{}ylabel}\PY{p}{(}\PY{l+s+s1}{\PYZsq{}}\PY{l+s+s1}{P\PYZhy{}Value}\PY{l+s+s1}{\PYZsq{}}\PY{p}{)}
\PY{n}{ax}\PY{o}{.}\PY{n}{set\PYZus{}xlabel}\PY{p}{(}\PY{l+s+s1}{\PYZsq{}}\PY{l+s+s1}{Index}\PY{l+s+s1}{\PYZsq{}}\PY{p}{)}
\PY{n}{ax}\PY{o}{.}\PY{n}{scatter}\PY{p}{(}\PY{n}{sorted\PYZus{}set\PYZus{}}\PY{o}{+}\PY{l+m+mi}{1}\PY{p}{,} \PY{n}{sorted\PYZus{}}\PY{p}{[}\PY{n}{sorted\PYZus{}set\PYZus{}}\PY{p}{]}\PY{p}{,} \PY{n}{c}\PY{o}{=}\PY{l+s+s1}{\PYZsq{}}\PY{l+s+s1}{r}\PY{l+s+s1}{\PYZsq{}}\PY{p}{,} \PY{n}{s}\PY{o}{=}\PY{l+m+mi}{20}\PY{p}{)}
\PY{n}{ax}\PY{o}{.}\PY{n}{axline}\PY{p}{(}\PY{p}{(}\PY{l+m+mi}{0}\PY{p}{,} \PY{l+m+mi}{0}\PY{p}{)}\PY{p}{,} \PY{p}{(}\PY{l+m+mi}{1}\PY{p}{,}\PY{n}{q}\PY{o}{/}\PY{n}{m}\PY{p}{)}\PY{p}{,} \PY{n}{c}\PY{o}{=}\PY{l+s+s1}{\PYZsq{}}\PY{l+s+s1}{k}\PY{l+s+s1}{\PYZsq{}}\PY{p}{,} \PY{n}{ls}\PY{o}{=}\PY{l+s+s1}{\PYZsq{}}\PY{l+s+s1}{\PYZhy{}\PYZhy{}}\PY{l+s+s1}{\PYZsq{}}\PY{p}{,} \PY{n}{linewidth}\PY{o}{=}\PY{l+m+mi}{3}\PY{p}{)}\PY{p}{;}
\end{Verbatim}
\end{tcolorbox}

    \begin{center}
    \adjustimage{max size={0.9\linewidth}{0.9\paperheight}}{artifacts/latex/Ch13-multiple-lab_files/artifacts/latex/Ch13-multiple-lab_46_0.png}
    \end{center}
    { \hspace*{\fill} \\}
    
    \hypertarget{a-re-sampling-approach}{%
\subsection{A Re-Sampling Approach}\label{a-re-sampling-approach}}

Here, we implement the re-sampling approach to hypothesis testing using
the \texttt{Khan} dataset, which we investigated in Section 13.5. First,
we merge the training and testing data, which results in observations on
83 patients for 2,308 genes.

    \begin{tcolorbox}[breakable, size=fbox, boxrule=1pt, pad at break*=1mm,colback=cellbackground, colframe=cellborder]
\prompt{In}{incolor}{23}{\boxspacing}
\begin{Verbatim}[commandchars=\\\{\}]
\PY{n}{Khan} \PY{o}{=} \PY{n}{load\PYZus{}data}\PY{p}{(}\PY{l+s+s1}{\PYZsq{}}\PY{l+s+s1}{Khan}\PY{l+s+s1}{\PYZsq{}}\PY{p}{)}      
\PY{n}{D} \PY{o}{=} \PY{n}{pd}\PY{o}{.}\PY{n}{concat}\PY{p}{(}\PY{p}{[}\PY{n}{Khan}\PY{p}{[}\PY{l+s+s1}{\PYZsq{}}\PY{l+s+s1}{xtrain}\PY{l+s+s1}{\PYZsq{}}\PY{p}{]}\PY{p}{,} \PY{n}{Khan}\PY{p}{[}\PY{l+s+s1}{\PYZsq{}}\PY{l+s+s1}{xtest}\PY{l+s+s1}{\PYZsq{}}\PY{p}{]}\PY{p}{]}\PY{p}{)}
\PY{n}{D}\PY{p}{[}\PY{l+s+s1}{\PYZsq{}}\PY{l+s+s1}{Y}\PY{l+s+s1}{\PYZsq{}}\PY{p}{]} \PY{o}{=} \PY{n}{pd}\PY{o}{.}\PY{n}{concat}\PY{p}{(}\PY{p}{[}\PY{n}{Khan}\PY{p}{[}\PY{l+s+s1}{\PYZsq{}}\PY{l+s+s1}{ytrain}\PY{l+s+s1}{\PYZsq{}}\PY{p}{]}\PY{p}{,} \PY{n}{Khan}\PY{p}{[}\PY{l+s+s1}{\PYZsq{}}\PY{l+s+s1}{ytest}\PY{l+s+s1}{\PYZsq{}}\PY{p}{]}\PY{p}{]}\PY{p}{)}
\PY{n}{D}\PY{p}{[}\PY{l+s+s1}{\PYZsq{}}\PY{l+s+s1}{Y}\PY{l+s+s1}{\PYZsq{}}\PY{p}{]}\PY{o}{.}\PY{n}{value\PYZus{}counts}\PY{p}{(}\PY{p}{)}
\end{Verbatim}
\end{tcolorbox}

            \begin{tcolorbox}[breakable, size=fbox, boxrule=.5pt, pad at break*=1mm, opacityfill=0]
\prompt{Out}{outcolor}{23}{\boxspacing}
\begin{Verbatim}[commandchars=\\\{\}]
Y
2    29
4    25
3    18
1    11
Name: count, dtype: int64
\end{Verbatim}
\end{tcolorbox}
        
    There are four classes of cancer. For each gene, we compare the mean
expression in the second class (rhabdomyosarcoma) to the mean expression
in the fourth class (Burkitt's lymphoma). Performing a standard
two-sample \(t\)-test\\
using \texttt{ttest\_ind()} from \texttt{scipy.stats} on the \(11\)th
gene produces a test-statistic of -2.09 and an associated \(p\)-value of
0.0412, suggesting modest evidence of a difference in mean expression
levels between the two cancer types.

    \begin{tcolorbox}[breakable, size=fbox, boxrule=1pt, pad at break*=1mm,colback=cellbackground, colframe=cellborder]
\prompt{In}{incolor}{24}{\boxspacing}
\begin{Verbatim}[commandchars=\\\{\}]
\PY{n}{D2} \PY{o}{=} \PY{n}{D}\PY{p}{[}\PY{k}{lambda} \PY{n}{df}\PY{p}{:}\PY{n}{df}\PY{p}{[}\PY{l+s+s1}{\PYZsq{}}\PY{l+s+s1}{Y}\PY{l+s+s1}{\PYZsq{}}\PY{p}{]} \PY{o}{==} \PY{l+m+mi}{2}\PY{p}{]}
\PY{n}{D4} \PY{o}{=} \PY{n}{D}\PY{p}{[}\PY{k}{lambda} \PY{n}{df}\PY{p}{:}\PY{n}{df}\PY{p}{[}\PY{l+s+s1}{\PYZsq{}}\PY{l+s+s1}{Y}\PY{l+s+s1}{\PYZsq{}}\PY{p}{]} \PY{o}{==} \PY{l+m+mi}{4}\PY{p}{]}
\PY{n}{gene\PYZus{}11} \PY{o}{=} \PY{l+s+s1}{\PYZsq{}}\PY{l+s+s1}{G0011}\PY{l+s+s1}{\PYZsq{}}
\PY{n}{observedT}\PY{p}{,} \PY{n}{pvalue} \PY{o}{=} \PY{n}{ttest\PYZus{}ind}\PY{p}{(}\PY{n}{D2}\PY{p}{[}\PY{n}{gene\PYZus{}11}\PY{p}{]}\PY{p}{,}
                              \PY{n}{D4}\PY{p}{[}\PY{n}{gene\PYZus{}11}\PY{p}{]}\PY{p}{,}
                              \PY{n}{equal\PYZus{}var}\PY{o}{=}\PY{k+kc}{True}\PY{p}{)}
\PY{n}{observedT}\PY{p}{,} \PY{n}{pvalue}
\end{Verbatim}
\end{tcolorbox}

            \begin{tcolorbox}[breakable, size=fbox, boxrule=.5pt, pad at break*=1mm, opacityfill=0]
\prompt{Out}{outcolor}{24}{\boxspacing}
\begin{Verbatim}[commandchars=\\\{\}]
(-2.0936330736768185, 0.04118643782678394)
\end{Verbatim}
\end{tcolorbox}
        
    However, this \(p\)-value relies on the assumption that under the null
hypothesis of no difference between the two groups, the test statistic
follows a \(t\)-distribution with \(29+25-2=52\) degrees of freedom.
Instead of using this theoretical null distribution, we can randomly
split the 54 patients into two groups of 29 and 25, and compute a new
test statistic. Under the null hypothesis of no difference between the
groups, this new test statistic should have the same distribution as our
original one. Repeating this process 10,000 times allows us to
approximate the null distribution of the test statistic. We compute the
fraction of the time that our observed test statistic exceeds the test
statistics obtained via re-sampling.

    \begin{tcolorbox}[breakable, size=fbox, boxrule=1pt, pad at break*=1mm,colback=cellbackground, colframe=cellborder]
\prompt{In}{incolor}{25}{\boxspacing}
\begin{Verbatim}[commandchars=\\\{\}]
\PY{n}{B} \PY{o}{=} \PY{l+m+mi}{10000}
\PY{n}{Tnull} \PY{o}{=} \PY{n}{np}\PY{o}{.}\PY{n}{empty}\PY{p}{(}\PY{n}{B}\PY{p}{)}
\PY{n}{D\PYZus{}} \PY{o}{=} \PY{n}{np}\PY{o}{.}\PY{n}{hstack}\PY{p}{(}\PY{p}{[}\PY{n}{D2}\PY{p}{[}\PY{n}{gene\PYZus{}11}\PY{p}{]}\PY{p}{,} \PY{n}{D4}\PY{p}{[}\PY{n}{gene\PYZus{}11}\PY{p}{]}\PY{p}{]}\PY{p}{)}
\PY{n}{n\PYZus{}} \PY{o}{=} \PY{n}{D2}\PY{p}{[}\PY{n}{gene\PYZus{}11}\PY{p}{]}\PY{o}{.}\PY{n}{shape}\PY{p}{[}\PY{l+m+mi}{0}\PY{p}{]}
\PY{n}{D\PYZus{}null} \PY{o}{=} \PY{n}{D\PYZus{}}\PY{o}{.}\PY{n}{copy}\PY{p}{(}\PY{p}{)}
\PY{k}{for} \PY{n}{b} \PY{o+ow}{in} \PY{n+nb}{range}\PY{p}{(}\PY{n}{B}\PY{p}{)}\PY{p}{:}
    \PY{n}{rng}\PY{o}{.}\PY{n}{shuffle}\PY{p}{(}\PY{n}{D\PYZus{}null}\PY{p}{)}
    \PY{n}{ttest\PYZus{}} \PY{o}{=} \PY{n}{ttest\PYZus{}ind}\PY{p}{(}\PY{n}{D\PYZus{}null}\PY{p}{[}\PY{p}{:}\PY{n}{n\PYZus{}}\PY{p}{]}\PY{p}{,}
                       \PY{n}{D\PYZus{}null}\PY{p}{[}\PY{n}{n\PYZus{}}\PY{p}{:}\PY{p}{]}\PY{p}{,}
                       \PY{n}{equal\PYZus{}var}\PY{o}{=}\PY{k+kc}{True}\PY{p}{)}
    \PY{n}{Tnull}\PY{p}{[}\PY{n}{b}\PY{p}{]} \PY{o}{=} \PY{n}{ttest\PYZus{}}\PY{o}{.}\PY{n}{statistic}
\PY{p}{(}\PY{n}{np}\PY{o}{.}\PY{n}{abs}\PY{p}{(}\PY{n}{Tnull}\PY{p}{)} \PY{o}{\PYZlt{}} \PY{n}{np}\PY{o}{.}\PY{n}{abs}\PY{p}{(}\PY{n}{observedT}\PY{p}{)}\PY{p}{)}\PY{o}{.}\PY{n}{mean}\PY{p}{(}\PY{p}{)}
\end{Verbatim}
\end{tcolorbox}

            \begin{tcolorbox}[breakable, size=fbox, boxrule=.5pt, pad at break*=1mm, opacityfill=0]
\prompt{Out}{outcolor}{25}{\boxspacing}
\begin{Verbatim}[commandchars=\\\{\}]
0.9602
\end{Verbatim}
\end{tcolorbox}
        
    This fraction, 0.0398, is our re-sampling-based \(p\)-value. It is
almost identical to the \(p\)-value of 0.0412 obtained using the
theoretical null distribution. We can plot a histogram of the
re-sampling-based test statistics in order to reproduce Figure 13.7.

    \begin{tcolorbox}[breakable, size=fbox, boxrule=1pt, pad at break*=1mm,colback=cellbackground, colframe=cellborder]
\prompt{In}{incolor}{26}{\boxspacing}
\begin{Verbatim}[commandchars=\\\{\}]
\PY{n}{fig}\PY{p}{,} \PY{n}{ax} \PY{o}{=} \PY{n}{plt}\PY{o}{.}\PY{n}{subplots}\PY{p}{(}\PY{n}{figsize}\PY{o}{=}\PY{p}{(}\PY{l+m+mi}{8}\PY{p}{,}\PY{l+m+mi}{8}\PY{p}{)}\PY{p}{)}
\PY{n}{ax}\PY{o}{.}\PY{n}{hist}\PY{p}{(}\PY{n}{Tnull}\PY{p}{,}
        \PY{n}{bins}\PY{o}{=}\PY{l+m+mi}{100}\PY{p}{,}
        \PY{n}{density}\PY{o}{=}\PY{k+kc}{True}\PY{p}{,}
        \PY{n}{facecolor}\PY{o}{=}\PY{l+s+s1}{\PYZsq{}}\PY{l+s+s1}{y}\PY{l+s+s1}{\PYZsq{}}\PY{p}{,}
        \PY{n}{label}\PY{o}{=}\PY{l+s+s1}{\PYZsq{}}\PY{l+s+s1}{Null}\PY{l+s+s1}{\PYZsq{}}\PY{p}{)}
\PY{n}{xval} \PY{o}{=} \PY{n}{np}\PY{o}{.}\PY{n}{linspace}\PY{p}{(}\PY{o}{\PYZhy{}}\PY{l+m+mf}{4.2}\PY{p}{,} \PY{l+m+mf}{4.2}\PY{p}{,} \PY{l+m+mi}{1001}\PY{p}{)}
\PY{n}{ax}\PY{o}{.}\PY{n}{plot}\PY{p}{(}\PY{n}{xval}\PY{p}{,}
        \PY{n}{t\PYZus{}dbn}\PY{o}{.}\PY{n}{pdf}\PY{p}{(}\PY{n}{xval}\PY{p}{,} \PY{n}{D\PYZus{}}\PY{o}{.}\PY{n}{shape}\PY{p}{[}\PY{l+m+mi}{0}\PY{p}{]}\PY{o}{\PYZhy{}}\PY{l+m+mi}{2}\PY{p}{)}\PY{p}{,}
        \PY{n}{c}\PY{o}{=}\PY{l+s+s1}{\PYZsq{}}\PY{l+s+s1}{r}\PY{l+s+s1}{\PYZsq{}}\PY{p}{)}
\PY{n}{ax}\PY{o}{.}\PY{n}{axvline}\PY{p}{(}\PY{n}{observedT}\PY{p}{,}
           \PY{n}{c}\PY{o}{=}\PY{l+s+s1}{\PYZsq{}}\PY{l+s+s1}{b}\PY{l+s+s1}{\PYZsq{}}\PY{p}{,}
           \PY{n}{label}\PY{o}{=}\PY{l+s+s1}{\PYZsq{}}\PY{l+s+s1}{Observed}\PY{l+s+s1}{\PYZsq{}}\PY{p}{)}
\PY{n}{ax}\PY{o}{.}\PY{n}{legend}\PY{p}{(}\PY{p}{)}
\PY{n}{ax}\PY{o}{.}\PY{n}{set\PYZus{}xlabel}\PY{p}{(}\PY{l+s+s2}{\PYZdq{}}\PY{l+s+s2}{Null Distribution of Test Statistic}\PY{l+s+s2}{\PYZdq{}}\PY{p}{)}\PY{p}{;}
\end{Verbatim}
\end{tcolorbox}

    \begin{center}
    \adjustimage{max size={0.9\linewidth}{0.9\paperheight}}{artifacts/latex/Ch13-multiple-lab_files/artifacts/latex/Ch13-multiple-lab_54_0.png}
    \end{center}
    { \hspace*{\fill} \\}
    
    The re-sampling-based null distribution is almost identical to the
theoretical null distribution, which is displayed in red.

Finally, we implement the plug-in re-sampling FDR approach outlined in
Algorithm 13.4. Depending on the speed of your computer, calculating the
FDR for all 2,308 genes in the \texttt{Khan} dataset may take a while.
Hence, we will illustrate the approach on a random subset of 100 genes.
For each gene, we first compute the observed test statistic, and then
produce 10,000 re-sampled test statistics. This may take a few minutes
to run. If you are in a rush, then you could set \texttt{B} equal to a
smaller value (e.g.~\texttt{B=500}).

    \begin{tcolorbox}[breakable, size=fbox, boxrule=1pt, pad at break*=1mm,colback=cellbackground, colframe=cellborder]
\prompt{In}{incolor}{27}{\boxspacing}
\begin{Verbatim}[commandchars=\\\{\}]
\PY{n}{m}\PY{p}{,} \PY{n}{B} \PY{o}{=} \PY{l+m+mi}{100}\PY{p}{,} \PY{l+m+mi}{10000}
\PY{n}{idx} \PY{o}{=} \PY{n}{rng}\PY{o}{.}\PY{n}{choice}\PY{p}{(}\PY{n}{Khan}\PY{p}{[}\PY{l+s+s1}{\PYZsq{}}\PY{l+s+s1}{xtest}\PY{l+s+s1}{\PYZsq{}}\PY{p}{]}\PY{o}{.}\PY{n}{columns}\PY{p}{,} \PY{n}{m}\PY{p}{,} \PY{n}{replace}\PY{o}{=}\PY{k+kc}{False}\PY{p}{)}
\PY{n}{T\PYZus{}vals} \PY{o}{=} \PY{n}{np}\PY{o}{.}\PY{n}{empty}\PY{p}{(}\PY{n}{m}\PY{p}{)}
\PY{n}{Tnull\PYZus{}vals} \PY{o}{=} \PY{n}{np}\PY{o}{.}\PY{n}{empty}\PY{p}{(}\PY{p}{(}\PY{n}{m}\PY{p}{,} \PY{n}{B}\PY{p}{)}\PY{p}{)}

\PY{k}{for} \PY{n}{j} \PY{o+ow}{in} \PY{n+nb}{range}\PY{p}{(}\PY{n}{m}\PY{p}{)}\PY{p}{:}
    \PY{n}{col} \PY{o}{=} \PY{n}{idx}\PY{p}{[}\PY{n}{j}\PY{p}{]}
    \PY{n}{T\PYZus{}vals}\PY{p}{[}\PY{n}{j}\PY{p}{]} \PY{o}{=} \PY{n}{ttest\PYZus{}ind}\PY{p}{(}\PY{n}{D2}\PY{p}{[}\PY{n}{col}\PY{p}{]}\PY{p}{,}
                          \PY{n}{D4}\PY{p}{[}\PY{n}{col}\PY{p}{]}\PY{p}{,}
                          \PY{n}{equal\PYZus{}var}\PY{o}{=}\PY{k+kc}{True}\PY{p}{)}\PY{o}{.}\PY{n}{statistic}
    \PY{n}{D\PYZus{}} \PY{o}{=} \PY{n}{np}\PY{o}{.}\PY{n}{hstack}\PY{p}{(}\PY{p}{[}\PY{n}{D2}\PY{p}{[}\PY{n}{col}\PY{p}{]}\PY{p}{,} \PY{n}{D4}\PY{p}{[}\PY{n}{col}\PY{p}{]}\PY{p}{]}\PY{p}{)}
    \PY{n}{D\PYZus{}null} \PY{o}{=} \PY{n}{D\PYZus{}}\PY{o}{.}\PY{n}{copy}\PY{p}{(}\PY{p}{)}
    \PY{k}{for} \PY{n}{b} \PY{o+ow}{in} \PY{n+nb}{range}\PY{p}{(}\PY{n}{B}\PY{p}{)}\PY{p}{:}
        \PY{n}{rng}\PY{o}{.}\PY{n}{shuffle}\PY{p}{(}\PY{n}{D\PYZus{}null}\PY{p}{)}
        \PY{n}{ttest\PYZus{}} \PY{o}{=} \PY{n}{ttest\PYZus{}ind}\PY{p}{(}\PY{n}{D\PYZus{}null}\PY{p}{[}\PY{p}{:}\PY{n}{n\PYZus{}}\PY{p}{]}\PY{p}{,}
                           \PY{n}{D\PYZus{}null}\PY{p}{[}\PY{n}{n\PYZus{}}\PY{p}{:}\PY{p}{]}\PY{p}{,}
                           \PY{n}{equal\PYZus{}var}\PY{o}{=}\PY{k+kc}{True}\PY{p}{)}
        \PY{n}{Tnull\PYZus{}vals}\PY{p}{[}\PY{n}{j}\PY{p}{,}\PY{n}{b}\PY{p}{]} \PY{o}{=} \PY{n}{ttest\PYZus{}}\PY{o}{.}\PY{n}{statistic}
\end{Verbatim}
\end{tcolorbox}

    Next, we compute the number of rejected null hypotheses \(R\), the
estimated number of false positives \(\widehat{V}\), and the estimated
FDR, for a range of threshold values \(c\) in Algorithm 13.4. The
threshold values are chosen using the absolute values of the test
statistics from the 100 genes.

    \begin{tcolorbox}[breakable, size=fbox, boxrule=1pt, pad at break*=1mm,colback=cellbackground, colframe=cellborder]
\prompt{In}{incolor}{28}{\boxspacing}
\begin{Verbatim}[commandchars=\\\{\}]
\PY{n}{cutoffs} \PY{o}{=} \PY{n}{np}\PY{o}{.}\PY{n}{sort}\PY{p}{(}\PY{n}{np}\PY{o}{.}\PY{n}{abs}\PY{p}{(}\PY{n}{T\PYZus{}vals}\PY{p}{)}\PY{p}{)}
\PY{n}{FDRs}\PY{p}{,} \PY{n}{Rs}\PY{p}{,} \PY{n}{Vs} \PY{o}{=} \PY{n}{np}\PY{o}{.}\PY{n}{empty}\PY{p}{(}\PY{p}{(}\PY{l+m+mi}{3}\PY{p}{,} \PY{n}{m}\PY{p}{)}\PY{p}{)}
\PY{k}{for} \PY{n}{j} \PY{o+ow}{in} \PY{n+nb}{range}\PY{p}{(}\PY{n}{m}\PY{p}{)}\PY{p}{:}
   \PY{n}{R} \PY{o}{=} \PY{n}{np}\PY{o}{.}\PY{n}{sum}\PY{p}{(}\PY{n}{np}\PY{o}{.}\PY{n}{abs}\PY{p}{(}\PY{n}{T\PYZus{}vals}\PY{p}{)} \PY{o}{\PYZgt{}}\PY{o}{=} \PY{n}{cutoffs}\PY{p}{[}\PY{n}{j}\PY{p}{]}\PY{p}{)}
   \PY{n}{V} \PY{o}{=} \PY{n}{np}\PY{o}{.}\PY{n}{sum}\PY{p}{(}\PY{n}{np}\PY{o}{.}\PY{n}{abs}\PY{p}{(}\PY{n}{Tnull\PYZus{}vals}\PY{p}{)} \PY{o}{\PYZgt{}}\PY{o}{=} \PY{n}{cutoffs}\PY{p}{[}\PY{n}{j}\PY{p}{]}\PY{p}{)} \PY{o}{/} \PY{n}{B}
   \PY{n}{Rs}\PY{p}{[}\PY{n}{j}\PY{p}{]} \PY{o}{=} \PY{n}{R}
   \PY{n}{Vs}\PY{p}{[}\PY{n}{j}\PY{p}{]} \PY{o}{=} \PY{n}{V}
   \PY{n}{FDRs}\PY{p}{[}\PY{n}{j}\PY{p}{]} \PY{o}{=} \PY{n}{V} \PY{o}{/} \PY{n}{R}
\end{Verbatim}
\end{tcolorbox}

    Now, for any given FDR, we can find the genes that will be rejected. For
example, with FDR controlled at 0.1, we reject 15 of the 100 null
hypotheses. On average, we would expect about one or two of these genes
(i.e.~10\% of 15) to be false discoveries. At an FDR of 0.2, we can
reject the null hypothesis for 28 genes, of which we expect around six
to be false discoveries.

The variable \texttt{idx} stores which genes were included in our 100
randomly-selected genes. Let's look at the genes whose estimated FDR is
less than 0.1.

    \begin{tcolorbox}[breakable, size=fbox, boxrule=1pt, pad at break*=1mm,colback=cellbackground, colframe=cellborder]
\prompt{In}{incolor}{29}{\boxspacing}
\begin{Verbatim}[commandchars=\\\{\}]
\PY{n+nb}{sorted}\PY{p}{(}\PY{n}{idx}\PY{p}{[}\PY{n}{np}\PY{o}{.}\PY{n}{abs}\PY{p}{(}\PY{n}{T\PYZus{}vals}\PY{p}{)} \PY{o}{\PYZgt{}}\PY{o}{=} \PY{n}{cutoffs}\PY{p}{[}\PY{n}{FDRs} \PY{o}{\PYZlt{}} \PY{l+m+mf}{0.1}\PY{p}{]}\PY{o}{.}\PY{n}{min}\PY{p}{(}\PY{p}{)}\PY{p}{]}\PY{p}{)}
\end{Verbatim}
\end{tcolorbox}

            \begin{tcolorbox}[breakable, size=fbox, boxrule=.5pt, pad at break*=1mm, opacityfill=0]
\prompt{Out}{outcolor}{29}{\boxspacing}
\begin{Verbatim}[commandchars=\\\{\}]
['G0097',
 'G0129',
 'G0182',
 'G0714',
 'G0812',
 'G0941',
 'G0982',
 'G1020',
 'G1022',
 'G1090',
 'G1320',
 'G1634',
 'G1697',
 'G1853',
 'G1854',
 'G1994',
 'G2017',
 'G2115',
 'G2193']
\end{Verbatim}
\end{tcolorbox}
        
    At an FDR threshold of 0.2, more genes are selected, at the cost of
having a higher expected proportion of false discoveries.

    \begin{tcolorbox}[breakable, size=fbox, boxrule=1pt, pad at break*=1mm,colback=cellbackground, colframe=cellborder]
\prompt{In}{incolor}{30}{\boxspacing}
\begin{Verbatim}[commandchars=\\\{\}]
\PY{n+nb}{sorted}\PY{p}{(}\PY{n}{idx}\PY{p}{[}\PY{n}{np}\PY{o}{.}\PY{n}{abs}\PY{p}{(}\PY{n}{T\PYZus{}vals}\PY{p}{)} \PY{o}{\PYZgt{}}\PY{o}{=} \PY{n}{cutoffs}\PY{p}{[}\PY{n}{FDRs} \PY{o}{\PYZlt{}} \PY{l+m+mf}{0.2}\PY{p}{]}\PY{o}{.}\PY{n}{min}\PY{p}{(}\PY{p}{)}\PY{p}{]}\PY{p}{)}
\end{Verbatim}
\end{tcolorbox}

            \begin{tcolorbox}[breakable, size=fbox, boxrule=.5pt, pad at break*=1mm, opacityfill=0]
\prompt{Out}{outcolor}{30}{\boxspacing}
\begin{Verbatim}[commandchars=\\\{\}]
['G0097',
 'G0129',
 'G0158',
 'G0182',
 'G0242',
 'G0552',
 'G0679',
 'G0714',
 'G0751',
 'G0812',
 'G0908',
 'G0941',
 'G0982',
 'G1020',
 'G1022',
 'G1090',
 'G1240',
 'G1244',
 'G1320',
 'G1381',
 'G1514',
 'G1634',
 'G1697',
 'G1768',
 'G1853',
 'G1854',
 'G1907',
 'G1994',
 'G2017',
 'G2115',
 'G2193']
\end{Verbatim}
\end{tcolorbox}
        
    The next line generates Figure 13.11, which is similar to Figure 13.9,
except that it is based on only a subset of the genes.

    \begin{tcolorbox}[breakable, size=fbox, boxrule=1pt, pad at break*=1mm,colback=cellbackground, colframe=cellborder]
\prompt{In}{incolor}{31}{\boxspacing}
\begin{Verbatim}[commandchars=\\\{\}]
\PY{n}{fig}\PY{p}{,} \PY{n}{ax} \PY{o}{=} \PY{n}{plt}\PY{o}{.}\PY{n}{subplots}\PY{p}{(}\PY{p}{)}
\PY{n}{ax}\PY{o}{.}\PY{n}{plot}\PY{p}{(}\PY{n}{Rs}\PY{p}{,} \PY{n}{FDRs}\PY{p}{,} \PY{l+s+s1}{\PYZsq{}}\PY{l+s+s1}{b}\PY{l+s+s1}{\PYZsq{}}\PY{p}{,} \PY{n}{linewidth}\PY{o}{=}\PY{l+m+mi}{3}\PY{p}{)}
\PY{n}{ax}\PY{o}{.}\PY{n}{set\PYZus{}xlabel}\PY{p}{(}\PY{l+s+s2}{\PYZdq{}}\PY{l+s+s2}{Number of Rejections}\PY{l+s+s2}{\PYZdq{}}\PY{p}{)}
\PY{n}{ax}\PY{o}{.}\PY{n}{set\PYZus{}ylabel}\PY{p}{(}\PY{l+s+s2}{\PYZdq{}}\PY{l+s+s2}{False Discovery Rate}\PY{l+s+s2}{\PYZdq{}}\PY{p}{)}\PY{p}{;}
\end{Verbatim}
\end{tcolorbox}

    \begin{center}
    \adjustimage{max size={0.9\linewidth}{0.9\paperheight}}{artifacts/latex/Ch13-multiple-lab_files/artifacts/latex/Ch13-multiple-lab_64_0.png}
    \end{center}
    { \hspace*{\fill} \\}
    
    


    % Add a bibliography block to the postdoc
    
    
    
\end{document}
